\documentclass[aps,prd,twocolumn,superscriptaddress,nofootinbib]{revtex4-2}

\usepackage{amsmath}
\usepackage{amssymb}
\usepackage{amsthm}
\usepackage{booktabs}
\usepackage{xcolor}
\usepackage{graphicx}
\usepackage{mathrsfs}

\newtheorem{theorem}{Theorem}[section]
\newtheorem{proposition}[theorem]{Proposition}
\newtheorem{lemma}[theorem]{Lemma}
\newtheorem{corollary}[theorem]{Corollary}
\newtheorem{conjecture}[theorem]{Conjecture}
\newtheorem{axiom}[theorem]{Axiom}
\theoremstyle{definition}
\newtheorem{definition}[theorem]{Definition}
\theoremstyle{remark}
\newtheorem{remark}[theorem]{Remark}
\newtheorem{example}[theorem]{Example}

\DeclareMathOperator{\ord}{ord}
\DeclareMathOperator{\rank}{rank}
\DeclareMathOperator{\Sel}{Sel}
\DeclareMathOperator{\Reg}{Reg}
\DeclareMathOperator{\Sha}{Sha}
\DeclareMathOperator{\Gal}{Gal}
\DeclareMathOperator{\Hom}{Hom}

\usepackage[hyperfootnotes=false]{hyperref}
\hypersetup{
    pdftitle={The BSD Conjecture as a Positivity Normal-Form: A Structural Survey},
    pdfauthor={Lukas Geiger},
    pdfsubject={Structural survey of the Birch and Swinnerton-Dyer conjecture via positivity normal forms},
    pdfkeywords={Birch and Swinnerton-Dyer conjecture, Neron-Tate height, positivity, Euler systems, Gross-Zagier, Iwasawa theory},
    colorlinks=true,
    linkcolor=black,
    urlcolor=blue,
    citecolor=black
}
\makeatletter
\renewcommand*{\HyperDestNameFilter}[1]{en.#1}
\makeatother

\begin{document}

% ===================================================================
% TITLE PAGE
% ===================================================================

\title{\texorpdfstring{The BSD Conjecture as a Positivity Normal-Form:\\A Structural Survey}{The BSD Conjecture as a Positivity Normal-Form: A Structural Survey}}

\author{Lukas Geiger}
\affiliation{Independent Researcher, Bernau, Germany\\ORCID: \href{https://orcid.org/0009-0005-7296-1534}{0009-0005-7296-1534}}

\date{\today}

\begin{abstract}
We reformulate the Birch and Swinnerton-Dyer conjecture as a \textit{positivity normal-form}: a structural characterisation, not a proof. Four axiomatic constraints -- finite arithmetic capacity, cooperative height positivity, controlled Iwasawa growth, and Euler-system composition coherence -- identify the BSD leading-term formula as the expected ``saturation profile'' once the missing bridge inputs are also assumed: height/cycle coupling, sufficiently strong Selmer/Euler-system control, and finiteness of $\Sha$ in general rank. For rank~$\geq 2$ these bridge inputs remain conjectural, and some axioms themselves overlap with BSD. The N\'eron-Tate height pairing provides the canonical positivity structure in this reading. We show that the Gross-Zagier theorem is the rank-1 instantiation of this framework, coupling an analytic derivative to a non-negative height. For general rank~$r$, we formulate the \textit{Height Saturation Conjecture}: $L^{(r)}(E,1)$ is proportional to the regulator $R_E = \det(\langle P_i, P_j \rangle) > 0$. This would force $\ord_{s=1} L(E,s) = \rank E(\mathbb{Q})$, but only after the missing higher Gross-Zagier and higher Kolyvagin/Sha inputs are supplied.
\newpage
\end{abstract}

\keywords{Birch and Swinnerton-Dyer conjecture, N\'eron-Tate height, positivity, Euler systems, Gross-Zagier, Iwasawa theory, normal-form characterisation}

\thanks{AI tools (Claude by Anthropic) were used for mathematical formalization, literature verification, and editorial refinement. All mathematical content was developed, verified, and approved by the author.}

\maketitle


% ===================================================================
% 1. INTRODUCTION
% ===================================================================
\section{Introduction}
\label{sec:intro}

\subsection{The BSD conjecture}

Let $E/\mathbb{Q}$ be an elliptic curve with conductor $N$ and $L$-function $L(E,s) = \sum_{n \geq 1} a_n n^{-s}$ (which exists for all $E/\mathbb{Q}$ by the modularity theorem \cite{Wiles1995}). The \textit{Birch and Swinnerton-Dyer conjecture} (BSD) \cite{BSD1965} asserts:

\begin{conjecture}[BSD -- Weak form]
\label{conj:bsd_weak}
$\ord_{s=1} L(E,s) = \rank E(\mathbb{Q})$.
\end{conjecture}

\begin{conjecture}[BSD -- Strong form]
\label{conj:bsd_strong}
The leading coefficient of $L(E,s)$ at $s = 1$ satisfies:
\begin{equation}
\frac{L^{(r)}(E,1)}{r!} = \frac{\Omega_E \cdot R_E \cdot \prod_p c_p \cdot |\Sha(E/\mathbb{Q})|}{|E(\mathbb{Q})_{\mathrm{tors}}|^2}\,,
\label{eq:bsd_formula}
\end{equation}
where $r = \rank E(\mathbb{Q})$, $\Omega_E$ is the real period, $R_E$ is the regulator (determinant of the N\'eron-Tate height pairing), $c_p$ are the Tamagawa numbers, and $\Sha(E/\mathbb{Q})$ is the Tate-Shafarevich group. We write $E_{\mathrm{tors}} := E(\mathbb{Q})_{\mathrm{tors}}$ throughout for brevity.
\end{conjecture}

The conjecture is known for:
\begin{itemize}
\item \textbf{Rank 0:} If $L(E,1) \neq 0$, then $\rank E(\mathbb{Q}) = 0$ and $|\Sha| < \infty$ (Kolyvagin \cite{Kolyvagin1990}, using Gross-Zagier \cite{GrossZagier1986}).
\item \textbf{Rank 1:} If $\ord_{s=1} L(E,s) = 1$, then $\rank E(\mathbb{Q}) = 1$ and $|\Sha| < \infty$ (Gross-Zagier + Kolyvagin).
\item \textbf{Rank $\geq 2$:} Open in general.
\end{itemize}

\subsection{The positivity perspective}

This paper proposes a structural lens in which the BSD conjecture, like parts of the Hodge and Weil stories, can be read through \textit{positivity}. The key observation:

\begin{quote}
\textit{The N\'eron-Tate height pairing is a canonical positive-definite quadratic form on the Mordell-Weil group. It detects arithmetic ``directions'' -- rational points modulo torsion -- via their positive energy cost. The BSD conjecture asserts that the $L$-function ``sees'' exactly these directions, no more and no fewer.}
\end{quote}

The parallel to existing positivity proofs is:

\begin{center}
\begin{tabular}{lll}
\toprule
Conjecture & Positive form & Result \\
\midrule
Weil & Frobenius eigenvalues & Deligne 1974 \\
Hodge & Hodge-Riemann form & Open \\
\textbf{BSD} & \textbf{N\'eron-Tate height} & \textbf{Rank $\leq 1$} \\
\bottomrule
\end{tabular}
\end{center}

The Gross-Zagier theorem \cite{GrossZagier1986} (extended to all Shimura curves by Yuan-Zhang-Zhang \cite{YZZ2013}) is the canonical example: it couples the first derivative $L'(E,1)$ to the height $\hat{h}(P_K)$ of a Heegner point, establishing:
\begin{equation}
L'(E,1) \neq 0 \iff \hat{h}(P_K) > 0 \iff P_K \text{ non-torsion}\,.
\end{equation}
This is a \textit{positivity identity}: an analytic object (derivative of $L$) equals a non-negative geometric object (height), and non-vanishing on one side forces non-triviality on the other.

\subsection{Normal-form theorems}

By a \textit{positivity normal-form theorem} we mean a result of the following type: given a set of structural axioms (positivity, finiteness, composition coherence), there exists a \textit{unique} functional relationship between the analytic and geometric invariants of the object in question, and this relationship is forced by the conjunction of the axioms. The Weil conjectures for the zeta function of a variety over $\mathbb{F}_q$ provide the archetype: Frobenius positivity (the Riemann hypothesis for varieties) forces the unique factorization of the zeta function via eigenvalues on \'etale cohomology.

\subsection{Structure of this paper}

We develop four axioms (Section~\ref{sec:axioms}) which, together with the bridge hypotheses made explicit later, select the BSD leading-term formula as the unique expected canonical-product normal form. We show how the rank $\leq 1$ case is already proven within this framework (Section~\ref{sec:rank01}), identify the precise obstructions for higher rank (Section~\ref{sec:higher_rank}), and discuss the relationship to Iwasawa theory and Euler systems (Section~\ref{sec:iwasawa}).

This paper is self-contained as a structural exposition: definitions and internal consistency checks are given here, while the genuine rank~$\leq 1$ inputs are the cited external theorems of Gross--Zagier, Kolyvagin, Kato and related Iwasawa-theoretic work. Connections to the broader FST programme are cited for context but are not logically required.


% ===================================================================
% 2. THE POSITIVITY STRUCTURE
% ===================================================================
\section{The Positivity Structure of BSD}
\label{sec:positivity}

\subsection{The N\'eron-Tate height}

Let $E/\mathbb{Q}$ be an elliptic curve. The \textit{N\'eron-Tate (canonical) height} $\hat{h}: E(\mathbb{Q}) \to \mathbb{R}_{\geq 0}$ is the unique function satisfying:
\begin{enumerate}
\item $\hat{h}(P) \geq 0$ for all $P \in E(\mathbb{Q})$, with equality iff $P$ is torsion.
\item $\hat{h}(nP) = n^2 \hat{h}(P)$ (quadratic scaling).
\item $\hat{h}$ differs from the naive (Weil) height by a bounded function.
\end{enumerate}

The associated \textit{N\'eron-Tate pairing} is:
\begin{equation}
\langle P, Q \rangle = \frac{1}{2}\left(\hat{h}(P+Q) - \hat{h}(P) - \hat{h}(Q)\right)\,.
\label{eq:NT_pairing}
\end{equation}
In particular, $\langle P, P \rangle = \hat{h}(P)$ for all $P \in E(\mathbb{Q})$.

\begin{theorem}[N\'eron-Tate positivity]
\label{thm:NT_positive}
The pairing $\langle \cdot, \cdot \rangle$ is positive semi-definite on $E(\mathbb{Q})$ and positive definite on $E(\mathbb{Q})/E(\mathbb{Q})_{\mathrm{tors}}$.
\end{theorem}

This is the BSD analog of the Hodge-Riemann bilinear form: a canonical, non-degenerate, positive-definite quadratic form on the ``interesting'' part of the arithmetic structure.

\subsection{The regulator as positivity detector}

Let $P_1, \ldots, P_r$ be generators of $E(\mathbb{Q})/E(\mathbb{Q})_{\mathrm{tors}}$, where $r = \rank E(\mathbb{Q})$. The \textit{regulator} is:
\begin{equation}
R_E = \det\!\left(\langle P_i, P_j \rangle\right)_{1 \leq i,j \leq r}\,.
\label{eq:regulator}
\end{equation}

\begin{proposition}[Regulator positivity]
\label{prop:reg_positive}
$R_E > 0$ whenever $r \geq 1$, and $R_E = 1$ by convention when $r = 0$.
\end{proposition}

\begin{proof}
The matrix $(\langle P_i, P_j \rangle)$ is the Gram matrix of a positive-definite form. Its determinant is strictly positive.
\end{proof}

The regulator is a \textit{positivity certificate} for the existence of $r$ independent rational points. It cannot vanish or be negative: positive-definiteness of $\hat{h}$ makes $R_E > 0$ an automatic consequence of $\rank E(\mathbb{Q}) = r$.

\subsection{The BSD formula as a positivity identity}

The strong BSD formula~\eqref{eq:bsd_formula} can be read as:
\begin{equation}
\underbrace{\frac{L^{(r)}(E,1)}{r!}}_{\text{analytic}} = \underbrace{\Omega_E}_{> 0} \cdot \underbrace{R_E}_{> 0} \cdot \underbrace{\prod c_p}_{\geq 1} \cdot \underbrace{\frac{|\Sha|}{|E_{\mathrm{tors}}|^2}}_{> 0}\,.
\label{eq:bsd_positivity}
\end{equation}

Every factor on the right is non-negative, and $R_E > 0$ whenever $r \geq 1$. The formula asserts that the analytic side -- which has no a priori reason to be non-negative or to have the ``right'' order of vanishing -- is \textit{forced} to match this positive structure.

This is precisely the pattern of a \textit{positivity normal-form theorem}: the analytic object is constrained by the positivity of the geometric/arithmetic object to which it is equal.


% ===================================================================
% 3. THE FOUR AXIOMS
% ===================================================================
\section{The Four Axioms}
\label{sec:axioms}

We formulate four structural conditions that any coherent arithmetic theory of $L$-functions and Mordell-Weil groups must satisfy.

\subsection{Axiom I: Finite Arithmetic Capacity}

\begin{axiom}[Finite Arithmetic Capacity]
\label{ax:finite}
The Tate-Shafarevich group $\Sha(E/\mathbb{Q})$ is finite:
\begin{equation}
|\Sha(E/\mathbb{Q})| < \infty\,.
\end{equation}
Equivalently in Selmer language, the exact sequence
\[
\begin{aligned}
0 &\longrightarrow E(\mathbb{Q})\otimes \mathbb{Q}_p/\mathbb{Z}_p \\
  &\longrightarrow \Sel_{p^\infty}(E/\mathbb{Q})
  \longrightarrow \Sha(E/\mathbb{Q})[p^\infty]
  \longrightarrow 0
\end{aligned}
\]
shows that finiteness of $\Sha$ implies
$\mathrm{corank}_{\mathbb{Z}_p}\Sel_{p^\infty}(E/\mathbb{Q})=\rank E(\mathbb{Q})$
and finite $p$-primary defect for every $p$. The corank equality alone is weaker;
the nontrivial content of Axiom~I is finiteness of the Tate--Shafarevich group.
(The $n$-Selmer groups $\Sel_n(E/\mathbb{Q})$ are always finite by a standard
argument using finiteness of the class group and the weak Mordell-Weil theorem.)
\end{axiom}

\textit{Status.} Finiteness of $\Sha$ is known for rank $\leq 1$ (Kolyvagin \cite{Kolyvagin1990}) and for certain higher-rank cases. It is expected to hold in general but remains unproven.

\textit{Role in the No-Go argument.} If $\Sha$ were infinite, one could imagine a curve with $\rank E(\mathbb{Q}) = 0$ but $\ord_{s=1} L(E,s) = 2$: the ``missing rank'' would be hidden in $\Sha$. Axiom~I excludes this loophole.

\subsection{Axiom II: Cooperative Height Positivity}

\begin{axiom}[Cooperative Height Positivity]
\label{ax:positivity}
The N\'eron-Tate height pairing is positive-definite on $E(\mathbb{Q})/E(\mathbb{Q})_{\mathrm{tors}}$:
\begin{equation}
\langle P, P \rangle > 0 \quad \forall\, P \in E(\mathbb{Q}) \setminus E(\mathbb{Q})_{\mathrm{tors}}\,.
\end{equation}
Moreover, the height pairing is ``cooperative'': heights of independent points combine quadratically via the Gram determinant $R_E = \det(\langle P_i, P_j \rangle) > 0$.
\end{axiom}

\textit{Status.} Unconditionally true (a theorem, not a conjecture).

\textit{Role in the No-Go argument.} Positivity ensures that the right-hand side of BSD has a definite sign structure: $R_E > 0$ is forced. Any analytic formula for $L^{(r)}(E,1)$ must be compatible with this positivity.

\subsection{Axiom III: Controlled Iwasawa Growth}

\begin{axiom}[Controlled Iwasawa Growth]
\label{ax:iwasawa}
In the cyclotomic $\mathbb{Z}_p$-extension $\mathbb{Q}_\infty/\mathbb{Q}$, the $p$-primary Selmer groups grow in a controlled, asymptotically predictable fashion. Writing $X_n = \Sel_{p^\infty}(E/\mathbb{Q}_n)^\vee$ for the Pontryagin dual (a finitely generated $\mathbb{Z}_p$-module whose $\mathbb{Z}_p$-rank equals $\rank E(\mathbb{Q}_n)$):
\begin{equation}
\mathrm{ord}_p\!\bigl(|X_n^{\mathrm{tors}}|\bigr) = \mu p^n + \lambda n + \nu + O(1)\,,
\label{eq:iwasawa_growth}
\end{equation}
where $X_n^{\mathrm{tors}}$ denotes the $\mathbb{Z}_p$-torsion submodule of $X_n$, and $\mu, \lambda \geq 0$ are the Iwasawa invariants. The Iwasawa main conjecture provides a $p$-adic $L$-function $L_p(E,s)$ whose characteristic ideal matches the Selmer module.
\end{axiom}

\textit{Status.} The Iwasawa main conjecture for elliptic curves has been established in important cases, with different results covering complementary aspects:
\begin{itemize}
\item \textbf{Kato \cite{Kato2004}:} Proves one divisibility direction ($\mathrm{char}_\Lambda(\Sel^\vee) \subseteq (\mathcal{L}_p)$, i.e., the characteristic ideal of the Selmer dual is contained in the ideal generated by the $p$-adic $L$-function) for all primes~$p$ at which the modular form is ordinary and the residual representation satisfies mild conditions. This holds without restriction on the rank and provides an upper bound on the Selmer module.
\item \textbf{Skinner-Urban \cite{SkinnerUrban2014}:} Prove the full main conjecture (both divisibilities, hence equality) for \textit{ordinary} primes~$p$ (i.e., $p \nmid a_p$), under technical conditions including that $\bar{\rho}_{E,p}$ is irreducible and that the residual representation satisfies certain local hypotheses. This covers a large class of elliptic curves at ordinary primes, but does not apply at supersingular primes.
\end{itemize}
For supersingular primes, the main conjecture requires a different formulation (Pollack, Kobayashi, Sprung) and is the subject of ongoing work. The full main conjecture for all primes and all elliptic curves remains open.

\textit{Role in the No-Go argument.} The controlled growth formula~\eqref{eq:iwasawa_growth} (note: ``controlled'' refers to the asymptotically predictable growth governed by $\mu$ and $\lambda$, not strict monotonicity, since the $O(1)$ error allows bounded fluctuations) constrains Selmer growth in the $\mathbb{Z}_p$-tower. Together with the relevant main conjecture, it couples the $p$-adic $L$-function to the Selmer module. For rank~$\geq 2$ this remains a $p$-adic input to the bridge problem, not by itself an archimedean rank-equality proof.

\subsection{Axiom IV: Euler-System Composition Coherence}

\begin{axiom}[Euler-System Composition Coherence]
\label{ax:euler}
In the cases where the bridge is available, one needs an Euler-system package $\{c_m\}_{m \in \mathcal{M}}$ -- a coherent family of cohomology classes indexed by moduli $m$ -- satisfying norm-compatibility relations:
\begin{equation}
\mathrm{cores}_{m\ell/m}(c_{m\ell}) = P_\ell(\mathrm{Frob}_\ell^{-1})\, c_m\,,
\label{eq:euler_norm}
\end{equation}
where $P_\ell$ is the Euler factor at $\ell$ and $\mathrm{cores}$ is the corestriction. When a sufficiently deep Kolyvagin derivative machinery is available, this package provides upper bounds on Selmer groups:
\begin{equation}
|\Sel(E/\mathbb{Q})| \leq C \cdot |\text{(image of Euler system)}|\,,
\end{equation}
for a computable constant $C$. This last bound is not contained in norm coherence alone; in higher rank it is precisely bridge hypothesis~(H2).
\end{axiom}

\textit{Status.} Euler systems exist for many elliptic curves:
\begin{itemize}
\item \textbf{Heegner points} (for rank 1): the classical Euler system \cite{Kolyvagin1990}.
\item \textbf{Kato's Euler system}: from $K$-theory of modular curves \cite{Kato2004}.
\item \textbf{Beilinson-Flach elements}: for Rankin-Selberg convolutions \cite{LLZ2014}.
\end{itemize}

\textit{Role in the No-Go argument.} Axiom~IV records the norm-coherent composition structure required of Euler-system classes. Where a sufficiently deep Kolyvagin derivative machinery is available, this structure gives upper bounds on Selmer groups. In rank~$\leq 1$ this is part of the established Kato/Gross--Zagier/Kolyvagin picture. In rank~$\geq 2$, norm coherence alone is not full Selmer control; the missing full-strength bound is precisely bridge hypothesis~(H2).

\begin{remark}[Root number and parity]
\label{rem:root_number}
The functional equation of $L(E,s)$ determines the \textit{root number} $w(E) = \pm 1$, which constrains the parity of $\ord_{s=1} L(E,s)$: if $w(E) = +1$, then $\ord_{s=1} L(E,s)$ is even; if $w(E) = -1$, it is odd. The \textit{parity conjecture} -- that $(-1)^{\rank E(\mathbb{Q})} = w(E)$ -- is proven in many cases (Nekov\'a\v{r} \cite{Nekovar2006}; T.~and V.~Dokchitser \cite{Dokchitser2010}). While this parity constraint is not one of our four axioms, it provides an additional consistency check: any ``alternative profile'' violating the parity conjecture would be excluded by the functional equation alone.
\end{remark}

\begin{remark}[Independence of the axioms]
\label{rem:axiom-independence}
For rank $\leq 1$, the BSD rank and finiteness conclusions are established by independent theorems (Gross--Zagier, Kolyvagin, Kato and related Euler-system results), and the axiom package can be instantiated in the usual Euler-system/Iwasawa settings. This does not mean that every formulation of Axiom~III is uniformly known for every prime and every curve. For rank $\geq 2$, Axioms~I and~IV are \textit{not} independent of BSD: finiteness of $\Sha$ (Axiom~I) is itself a consequence of the full BSD conjecture, and the existence of a complete Kolyvagin system of sufficient depth (Axiom~IV) is not established. The conditional result (Proposition~\ref{thm:bsd_normal}) should therefore be read as a \textit{structural characterisation}: it identifies the BSD product formula as the expected normal form within the stated bridge package, rather than providing an independent proof of BSD.
\end{remark}

\subsection{Heuristic parallels}
\label{sec:heuristic}

The four axioms exhibit structural parallels to the variational framework of~\cite{FST-Unified2026}. These parallels are \textit{purely heuristic and motivational} -- they suggest why the axiom structure is natural, but do not constitute rigorous mathematical connections and are not used in any proof.

\begin{itemize}
\item \textbf{Axiom~I $\leftrightarrow$ Finite geometric capacity (Axiom~A).} The finiteness of $\Sha$ is the arithmetic analog of finite capacity: the ``hidden mass'' is bounded, preventing infinite discrepancies from being absorbed into the right-hand side of the BSD formula.

\item \textbf{Axiom~II $\leftrightarrow$ Cooperative reinforcement (Axiom~B).} Each independent rational point costs positive energy (height). The quadratic scaling $\hat{h}(nP) = n^2 \hat{h}(P)$ mirrors cooperative growth in the saturation ODE.

\item \textbf{Axiom~III $\leftrightarrow$ Controlled growth (Axiom~C).} As one ascends the $\mathbb{Z}_p$-tower, the torsion and corank data of Selmer modules are constrained by Iwasawa invariants. This is controlled arithmetic growth, not a direct prohibition of all rank jumps.

\item \textbf{Axiom~IV $\leftrightarrow$ Composition law (Axiom~D).} Norm-compatible steps in the Euler system compose coherently. The norm-compatibility~\eqref{eq:euler_norm} ensures that the Euler system is a globally consistent family, not a collection of ad hoc local constructions.
\end{itemize}


% ===================================================================
% 4. THE NORMAL-FORM THEOREM
% ===================================================================
\section{The Normal-Form Theorem}
\label{sec:normal_form}

\subsection{Statement}

\begin{proposition}[BSD Normal-Form Compatibility -- Conditional Template]
\label{thm:bsd_normal}
\textbf{Separate the assumptions into two layers.} Axioms~I--IV specify the
finiteness, positivity, Iwasawa-control, and Euler-system coherence constraints
that any BSD-type leading-term formula must respect. The passage from these
constraints to the BSD leading term additionally requires bridge hypotheses:
(H1) a Gross--Zagier-type height coupling, (H2) higher-rank
Kolyvagin/Selmer-control machinery, and (H3) finiteness of $\Sha$ in general
rank. For rank~$\leq 1$, the required bridge input is supplied by the
established Gross--Zagier, Kato, and Kolyvagin machinery. For rank~$\geq 2$,
if Axioms~I--IV and the bridge hypotheses (H1)--(H3) hold, then the BSD
leading-term expression is the canonical normal form compatible with these
constraints. Without (H1)--(H3), the result is a structural compatibility
diagnosis, not a proof of BSD or of rank equality.

Let $E/\mathbb{Q}$ be an elliptic curve satisfying Axioms~I--IV and the bridge hypotheses (H1)--(H3). Then:
\begin{enumerate}
\item The analytic rank equals the algebraic rank:
\begin{equation}
\ord_{s=1} L(E,s) = \rank E(\mathbb{Q})\,.
\end{equation}
\item The leading term has the normal form:
\begin{equation}
\frac{L^{(r)}(E,1)}{r!} = \Omega_E \cdot R_E \cdot \prod_p c_p \cdot \frac{|\Sha(E/\mathbb{Q})|}{|E(\mathbb{Q})_{\mathrm{tors}}|^2}\,,
\end{equation}
with all factors on the right determined by the arithmetic of $E$.
\item No alternative leading-term profile of the canonical BSD-product type is compatible with Axioms~I--IV plus (H1)--(H3).
\end{enumerate}
\end{proposition}

\begin{remark}[Conditional status and logical circularity at rank~$\geq 2$]
\label{rem:conditional_status}
Proposition~\ref{thm:bsd_normal} is a \textit{conditional structural characterisation}: it shows that the BSD formula is the expected normal form \textit{given} Axioms~I--IV and the bridge hypotheses (H1)--(H3). \textbf{It does not provide an unconditional proof of BSD, and for rank~$\geq 2$ the argument is logically circular if (H1)--(H3) are treated as already available.} For rank~$\leq 1$, the corresponding bridge is supplied by Gross--Zagier and Kolyvagin/Kato-type theorems, so the conclusion is a genuine structural reformulation of known results. For rank~$\geq 2$, however, Axiom~I/H3 (finiteness of $\Sha$) is itself a consequence of the full BSD conjecture, and Axiom~IV/H2 (existence of a complete Kolyvagin system of sufficient depth) is not established independently of BSD. The proposition therefore has the logical structure ``BSD-level bridge inputs + known facts $\Rightarrow$ BSD has the expected product normal form'' at rank~$\geq 2$, which is logically valid but mathematically trivial as a proof strategy. Its value lies in \textit{structural analysis}: identifying the exact missing bridge inputs, even though it does not provide them.

The analogy to the Weil conjectures requires qualification: Grothendieck's reformulation via \'etale cohomology was \textit{not} circular -- \'etale cohomology was an independently constructed mathematical object, and the reformulation proved rationality and the functional equation outright, reducing the problem to the Riemann hypothesis for varieties. By contrast, our axioms are not independently established for rank~$\geq 2$. The analogy holds \textit{in spirit} (identifying the correct structural framework) but not \textit{in logical force}.
\end{remark}

\subsection{Proof structure}

The proof strategy has three components, mirroring the saturation theorem:

\textbf{Step 1: Positivity forces the sign.} By Axiom~II, $R_E > 0$ whenever $r \geq 1$. By Axiom~I, $|\Sha| < \infty$, so the right-hand side of the BSD formula is a positive real number (for $r \geq 1$) or equals $L(E,1)/\Omega_E$ times arithmetic factors (for $r = 0$). The analytic side must match this sign.

\textbf{Step 2: Composition forces the order.} \textit{[Status: open for rank $\geq 2$; proven for rank $\leq 1$.]} For rank~0, Kato's Euler system \cite{Kato2004} shows that $L(E,1) \neq 0$ implies $\rank E(\mathbb{Q}) = 0$. For rank~1, Kolyvagin's Euler system \cite{Kolyvagin1990} shows that if $\ord_{s=1} L(E,s) = 1$ (and the Heegner point is non-torsion), then $\rank E(\mathbb{Q}) = 1$ and $|\Sha| < \infty$. \textbf{For rank~$\geq 2$, the general inequality $\rank E(\mathbb{Q}) \leq \ord_{s=1} L(E,s)$ is not established} -- it is equivalent to one direction of the BSD conjecture and requires the unproven hypothesis~(H2). The reverse inequality $\ord_{s=1} L(E,s) \leq \rank E(\mathbb{Q})$ requires the height coupling (H1): if $L^{(r)}(E,1)$ is proportional to $R_E > 0$ (Axiom~II), then $L$ cannot vanish to order higher than $r = \rank E(\mathbb{Q})$ without contradicting the positivity of the regulator. \textit{Note: both directions of this step are conditional on unproven hypotheses for rank~$\geq 2$ (see Remark below).}

\begin{remark}[Heuristic nature of the reverse inequality in Step~2]
The direction $\ord_{s=1} L(E,s) \leq \rank E(\mathbb{Q})$ relies on the full BSD leading-term formula as an assumption: without the BSD formula, the argument that ``the leading term vanishes faster than the regulator, violating positivity'' remains heuristic. Step~2 should therefore be understood as conditional on the BSD formula, not as a self-contained deduction from Axioms~I--IV alone (cf.\ Remark~\ref{rem:conditional_status} for the full discussion of logical dependencies).
\end{remark}

\textbf{Step 3: Uniqueness of the normal form.} Given the rank equality, the BSD formula is the unique expression compatible with:
\begin{itemize}
\item The functional equation and period normalisation of $L(E,s)$ (fix the archimedean period factor $\Omega_E$ once the BSD-type leading-term package is assumed).
\item The Birch lemma and its generalizations (determines the local factors $c_p$).
\item The regulator $R_E$ (determined by the height pairing).
\item Finiteness of $\Sha$ (Axiom~I).
\end{itemize}

\subsection{Uniqueness of the leading-term profile}

\begin{proposition}[Uniqueness of the leading-term formula -- Proof sketch]
\label{thm:uniqueness}
Let $E/\mathbb{Q}$ be an elliptic curve with $\rank E(\mathbb{Q}) = r$, satisfying Axioms~I--IV and the height coupling hypothesis~(H1). Then the BSD leading-term formula~\eqref{eq:bsd_formula} is the \textit{unique} expression
\[
  \frac{L^{(r)}(E,1)}{r!} = \Phi(E)
\]
satisfying the following five constraints simultaneously:
\begin{enumerate}
  \item[\textbf{(U1)}] \textbf{Analytic-algebraic rank equality:} $\ord_{s=1} L(E,s) = \rank E(\mathbb{Q})$.
  \item[\textbf{(U2)}] \textbf{Positivity of leading coefficient:} $\Phi(E) > 0$ whenever $r \geq 1$ (forced by Axiom~II, since any coupling to $R_E > 0$ must preserve the sign).
  \item[\textbf{(U3)}] \textbf{Integrality of the $\Sha$ contribution:} the ratio $\Phi(E)/(\Omega_E \cdot R_E \cdot \prod_p c_p)$ is of the form $|\Sha|/|E(\mathbb{Q})_{\mathrm{tors}}|^2$ with $|\Sha| \in \mathbb{Z}_{>0}$ (forced by Axiom~I).
  \item[\textbf{(U4)}] \textbf{Compatibility with functional equation:} $\Phi(E)$ is compatible with the $p$-adic interpolation formula $L_p(E,0) \sim L(E,1)/\Omega_E^+$ and its higher derivatives (forced by Axiom~III, the Iwasawa main conjecture specialization).
  \item[\textbf{(U5)}] \textbf{Isogeny invariance:} $\Phi(E)$ depends only on the isogeny class of~$E$, not on the choice of Weierstrass model or basis of $E(\mathbb{Q})$.
\end{enumerate}
\end{proposition}

\begin{proof}[Heuristic argument (not a rigorous proof)]
We argue \textit{heuristically} that (U1)--(U5) jointly determine $\Phi(E)$. \textbf{Caveat:} This argument is a plausibility sketch, not a rigorous proof. The rigorous framework for the uniqueness question is the \textit{Bloch-Kato conjecture} (Tamagawa number conjecture) \cite{BlochKato1990,FontainePerrinRiou1994}, which defines Tamagawa numbers for motives and derives the BSD formula as the special case for the motive $h^1(E)(1)$, eliminating the rational ambiguity of the Beilinson conjectures. Our heuristic argument should be read as motivation for \textit{why} the BSD formula is the expected answer; the Bloch-Kato framework provides the correct mathematical setting for making this precise.

\textit{Step~1: The positivity factor must be $R_E$.} By (U2), $\Phi(E) > 0$ for $r \geq 1$, so $\Phi(E)$ must involve a positive-definite invariant of $E(\mathbb{Q})/E(\mathbb{Q})_{\mathrm{tors}}$. The only canonical such invariant is the Gram determinant $R_E = \det(\langle P_i, P_j \rangle)$ of the N\'eron-Tate height, which is the \textit{unique} positive-definite bilinear form on $E(\mathbb{Q})/E(\mathbb{Q})_{\mathrm{tors}}$ satisfying quadratic scaling $\hat{h}(nP) = n^2 \hat{h}(P)$, agreement with the Weil height up to $O(1)$, and positive-definiteness modulo torsion (cf.\ \cite{Silverman2009}, Theorem~VIII.9.3). \textit{Note:} This step only shows that \textit{if} $\Phi(E)$ contains a positive-definite factor, it must be $R_E$; it does not show that $\Phi(E)$ must contain such a factor at all. The latter requires the height coupling hypothesis~(H1).

\textit{Step~2: The archimedean and local factors are determined.} (U4) (compatibility with the $p$-adic interpolation formula and functional equation) determines the archimedean factor $\Omega_E$ and the local factors $c_p$ through the Birch lemma and the specialization of the Iwasawa main conjecture.

\textit{Step~3: The remaining factor is constrained by integrality.} After extracting $\Omega_E \cdot R_E \cdot \prod_p c_p$, (U3) requires the remaining ratio $\Phi(E)/(\Omega_E \cdot R_E \cdot \prod_p c_p)$ to equal $|\Sha(E/\mathbb{Q})|/|E(\mathbb{Q})_{\mathrm{tors}}|^2$ with $|\Sha| \in \mathbb{Z}_{>0}$. \textit{Note:} This step presupposes that the decomposition $\Phi(E) = \Omega_E \cdot R_E \cdot (\text{rest})$ is valid, which is precisely the content of the BSD formula. The ``integrality constraint'' does not independently derive this decomposition; rather, it is \textit{consistent with} the BSD formula once assumed.

\textit{Step~4: Isogeny invariance excludes ad hoc constructions.} (U5) ensures that $\Phi(E)$ depends only on the isogeny class, not on auxiliary choices (basis of $E(\mathbb{Q})$, Weierstrass model). The individual quantities $\Omega_E$, $R_E$, $c_p$, $|\Sha|$, $|E(\mathbb{Q})_{\mathrm{tors}}|$ are \textit{not} separately isogeny-invariant, but their combination in the BSD formula is, since $L(E,s)$ depends only on the isogeny class (Faltings). This forces $\Phi(E)$ to be the unique isogeny-invariant product of canonical arithmetic data.

In summary, the five constraints make the BSD formula the unique \textit{natural candidate} of the form ``product of canonical arithmetic invariants.'' This is a weaker claim than a rigorous uniqueness proof: it shows that the BSD formula is the unique expression \textit{within the class of products of known canonical invariants}, not that it is the unique formula satisfying (U1)--(U5) in the space of all possible functions of~$E$. The Bloch-Kato conjecture \cite{BlochKato1990} provides the rigorous motivic framework in which this uniqueness can be made precise.
\end{proof}

\subsection{Exclusion of alternatives}

We show that ``alternative profiles'' -- hypothetical scenarios where the BSD formula fails -- violate the axioms:

\begin{proposition}[Consistency check: scenarios excluded by the axioms]
\label{prop:exclusion}
The following scenarios are excluded by the axioms. \textit{Note:} For rank~$\leq 1$, the exclusions restate known results (Kato, Kolyvagin) in the language of our axiom system. For rank~$\geq 2$, only (E3) and (E4) are excluded, and both follow trivially from the axiom assumptions. The non-trivial cases (E1) and (E2) remain conditional for rank~$\geq 2$.

\begin{enumerate}
\item[\textbf{(E1)}] \textbf{Analytic zero without geometric direction:} $\ord_{s=1} L(E,s) = r > \rank E(\mathbb{Q})$. Note that Euler-system bounds (Axiom~IV) establish the \textit{opposite} inequality $\rank E(\mathbb{Q}) \leq \ord_{s=1} L(E,s)$, so they do not directly exclude this scenario. Rather, exclusion of (E1) requires Axiom~I ($|\Sha| < \infty$) combined with the height coupling hypothesis~(H1): if the coupling $L^{(r)}(E,1) = C_r \cdot R_E$ holds (where $r = \rank E(\mathbb{Q})$), then $\ord > r$ would force $L^{(r)}(E,1) = 0$ (since $L$ vanishes to order $> r$), while $R_E > 0$ (Axiom~II) and $C_r > 0$, yielding a contradiction. For rank $\leq 1$, Gross-Zagier provides this coupling explicitly. \textit{Note:} For rank $\geq 2$, this exclusion is conditional on (H1) and (H2).

\item[\textbf{(E2)}] \textbf{Geometric rank without analytic derivative:} $\rank E(\mathbb{Q}) = r > \ord_{s=1} L(E,s) = s$. \textit{Under the assumption of the leading-term formula} and the Euler-system bound (H2), this is excluded: (H2) gives $r = \rank E(\mathbb{Q}) \leq \ord_{s=1} L(E,s) = s$, directly contradicting $r > s$. Thus for rank~$\leq 1$ (where Axiom~IV is proven unconditionally via Kato's Euler system), (E2) is excluded outright. \textbf{For rank~$\geq 2$, exclusion is conditional:} it depends on the full Euler-system bound (H2), which itself requires the existence of a Kolyvagin system of sufficient depth -- a major open problem. Without (H2), the inequality $\rank \leq \ord$ is not established, and no contradiction arises; thus (E2) \textit{cannot currently be excluded} for rank~$\geq 2$. Supplementarily, under the assumption that the leading-term formula holds, the height coupling (H1) provides a compatible perspective: it asserts that $L^{(r)}(E,1)/r!$ equals $C_r \cdot R_E$ as a \textit{leading-term identity}, coupling the $r$-th derivative to the regulator precisely when $\ord_{s=1} L(E,s) = r$; if $\ord = s < r$, the hypothesis of the coupling formula is not met.

\item[\textbf{(E3)}] \textbf{Infinite $\Sha$ absorbing discrepancy:} $\ord_{s=1} L(E,s) = r$ but $|\Sha| = \infty$, allowing a different right-hand side. Excluded by Axiom~I.

\item[\textbf{(E4)}] \textbf{Oscillating rank in $\mathbb{Z}_p$-towers:} Rank jumps non-monotonically through cyclotomic layers. Excluded by Axiom~III.
\end{enumerate}
\end{proposition}


% ===================================================================
% 5. RANK 0 AND 1: THE PROVEN CASES
% ===================================================================
\section{\texorpdfstring{Rank $\leq 1$: Positivity in Action}{Rank <= 1: Positivity in Action}}
\label{sec:rank01}

\subsection{Rank 0: Kato's exclusion}

\begin{proposition}[Rank 0 saturation]
\label{prop:rank0}
Let $E/\mathbb{Q}$ with $L(E,1) \neq 0$. Then all four axioms are verified, $\rank E(\mathbb{Q}) = 0$, $|\Sha| < \infty$, and the BSD formula holds.
\end{proposition}

\begin{proof}[Proof sketch]
\textit{Axiom II applied:} If $\rank E(\mathbb{Q}) \geq 1$, there exists a non-torsion point $P$ with $\hat{h}(P) > 0$. However, the direct implication $L(E,1) \neq 0 \Rightarrow \rank = 0$ is not established via Axiom~II alone; rather, it follows from the Euler-system machinery (Axiom~IV) below.

\textit{Axiom IV applied:} Kato's Euler system \cite{Kato2004}, constructed from $K$-theory classes on modular curves, provides the relevant upper bound for rank~0. Specifically, for any prime~$p$ at which the residual representation $\bar{\rho}_{E,p}$ is surjective (which holds for all but finitely many~$p$), Kato shows that if $L(E,1) \neq 0$, the image of his Euler system class in $H^1(\mathbb{Q}, T_p(E))$ is non-trivial, which forces the $p^\infty$-Selmer group to have $\mathbb{Z}_p$-corank zero:
\begin{equation}
\mathrm{corank}_{\mathbb{Z}_p}\, \Sel_{p^\infty}(E/\mathbb{Q}) = 0\,,
\end{equation}
hence $\rank E(\mathbb{Q}) = 0$. Combined with the Selmer bound, this also yields $|\Sha(E/\mathbb{Q})| < \infty$. (Note: for rank~0, the Heegner-point Euler system is not directly applicable since the Heegner point is torsion when $L(E,1) \neq 0$; Kato's Euler system is the appropriate tool.)

The rank-0 case is fully ``saturated'' in the positivity framework: all four axioms are verified (Axiom~III at all primes of good ordinary reduction with surjective residual representation; the remaining primes are handled by complementary results or are finite in number), and the No-Go argument excludes $\rank \geq 1$.
\end{proof}

\subsection{Rank 1: Gross-Zagier as a positivity identity}

The Gross-Zagier theorem \cite{GrossZagier1986} states: for a suitable imaginary quadratic field $K$ satisfying the Heegner hypothesis (every prime dividing the conductor $N$ splits in $K$ -- such a $K$ always exists by the Chinese Remainder Theorem and quadratic reciprocity),
\begin{equation}
L'(E,1) = \frac{8\pi^2 \|f\|^2}{\sqrt{|D_K|}} \cdot \hat{h}(P_K)\,,
\label{eq:gross_zagier}
\end{equation}
where $f$ is the newform associated to $E$, $\|f\|^2 = \int_{\Gamma_0(N) \backslash \mathcal{H}} |f(z)|^2 \, \frac{dx\,dy}{y^2}$ is the Petersson norm (the $L^2$-norm of $f$ with respect to the hyperbolic measure on the modular curve), $D_K$ is the discriminant, and $P_K \in E(K)$ is the Heegner point. (The formula is stated up to an explicit nonzero rational factor depending on $K$; for $|D_K| > 4$ this factor equals~$1$. The relationship between the Petersson norm $\|f\|^2$ and the N\'eron period $\Omega_E$ in the BSD formula involves the Manin constant~$c_M$; for the optimal curve in each isogeny class, $c_M = 1$ is verified computationally for all conductors up to $500{,}000$ in the Cremona database \cite{LMFDB}. Cremona's \texttt{ecdata} provides complete tables up to conductor $N = 500{,}000$; see \url{https://johncremona.github.io/ecdata/}. This verification is unconditional for $N \leq 130{,}000$ and conditional on optimality of the first curve for $N > 400{,}000$. More generally, \v{C}esnavi\v{c}ius~\cite{Cesnavicius2018} proved $c_M = 1$ for all semistable elliptic curves over~$\mathbb{Q}$.)

\begin{proposition}[Gross-Zagier as positivity identity]
\label{prop:gz_positivity}
The Gross-Zagier formula~\eqref{eq:gross_zagier} is a positivity identity:
\begin{itemize}
\item The left side $L'(E,1)$ has no a priori sign constraint.
\item The right side is a product of $\hat{h}(P_K) \geq 0$ (by Axiom~II) and positive constants.
\item Non-vanishing: $L'(E,1) \neq 0 \iff \hat{h}(P_K) > 0 \iff P_K$ is non-torsion.
\end{itemize}
\end{proposition}

This is the BSD analog of the Hodge-Riemann bilinear relation: the analytic object is equal to a geometrically positive quantity, and the positivity forces the correct rank.

Kolyvagin's Euler system then completes the argument: it shows $\rank E(\mathbb{Q}) = 1$ (the Heegner point generates a finite-index subgroup) and $|\Sha| < \infty$.

\begin{remark}
The rank-1 proof follows the exact pattern of a No-Go theorem:
\begin{enumerate}
\item Assume $\ord_{s=1} L(E,s) = 1$.
\item Gross-Zagier converts the analytic datum to a positive geometric datum ($\hat{h}(P_K) > 0$).
\item Kolyvagin excludes alternative explanations ($\rank > 1$ or $|\Sha| = \infty$).
\item Conclusion: $\rank E(\mathbb{Q}) = 1$, and the BSD formula holds.
\end{enumerate}
No explicit construction of rational points is needed beyond the Heegner point (which serves as a ``seed'' for the Euler system, not as a constructive proof of rank).
\end{remark}

\subsection{Numerical illustration}

\begin{example}[Curve 11a -- rank 0]
\label{ex:11a}
The elliptic curve $E: y^2 + y = x^3 - x^2 - 10x - 20$ (Cremona label 11a1, conductor $N = 11$; \cite{LMFDB}) is the curve of smallest conductor. It has $\rank E(\mathbb{Q}) = 0$, $E(\mathbb{Q})_{\mathrm{tors}} \cong \mathbb{Z}/5\mathbb{Z}$, and $|\Sha| = 1$ (proven). With $R_E = 1$ (by convention for rank~0), $\Omega_E = 1.2692\ldots$, $c_{11} = 5$ (Kodaira type $I_5$, split multiplicative reduction), and $|E(\mathbb{Q})_{\mathrm{tors}}| = 5$:
\[
\begin{aligned}
  L(E,1)
  &= \Omega_E \cdot R_E \cdot c_{11}
     \cdot \frac{|\Sha|}{|E_{\mathrm{tors}}|^2} \\
  &= 1.2692 \cdot 1 \cdot 5 \cdot \frac{1}{25}
   = 0.2538\ldots
\end{aligned}
\]
This matches the numerical value $L(E,1) = 0.2538\ldots$ to all available precision. All four axioms are verified: $|\Sha| = 1$ (I), height positivity vacuously satisfied for $r = 0$ (II), Iwasawa main conjecture proven at all primes $p \neq 11$ (III), Kato's Euler system provides the Selmer bound (IV).

The positivity interpretation: rank~0 is the ``vacuum state'' of the positivity framework -- no height direction exists, $R_E = 1$ by convention, and the BSD formula reduces to a ratio of archimedean and local data. The positivity structure is trivially saturated.
\end{example}

\begin{example}[Curve 37a -- rank 1]
\label{ex:37a}
The elliptic curve $E: y^2 + y = x^3 - x$ (Cremona label 37a1, conductor $N = 37$) has $\rank E(\mathbb{Q}) = 1$, generated by $P = (0,0)$ with $\hat{h}(P) = 0.0511\ldots$ The BSD data: $\Omega_E = 5.9869\ldots$, $R_E = 0.0511\ldots$, $c_{37} = 1$, $|E(\mathbb{Q})_{\mathrm{tors}}| = 1$, $|\Sha| = 1$ (proven). The BSD formula predicts $L'(E,1) = \Omega_E \cdot R_E \cdot 1 \cdot 1/1 = 0.3059\ldots$, which matches the numerical computation $L'(E,1) = 0.3059\ldots$ to all available digits. All four axioms are verified: $|\Sha| = 1 < \infty$ (I), $\hat{h}(P) > 0$ (II), Iwasawa main conjecture proven for $p \neq 37$ (III), Heegner point Euler system (IV).
\end{example}

\begin{example}[Curve 389a -- rank 2]
\label{ex:389a}
The elliptic curve $E: y^2 + y = x^3 + x^2 - 2x$ (Cremona label 389a1) has $\rank E(\mathbb{Q}) = 2$, generated by $P_1 = (-1,1)$ and $P_2 = (0,0)$. The height matrix depends on the choice of generators; the regulator $R_E$ is invariant under change of $\mathbb{Z}$-basis of $E(\mathbb{Q})/E(\mathbb{Q})_{\mathrm{tors}}$ because $\det(A^T M A) = (\det A)^2 \det M$ and $\det A = \pm 1$ for $A \in \mathrm{GL}(r, \mathbb{Z})$. The height matrix is:
\[
\left(\langle P_i, P_j \rangle\right) = \begin{pmatrix} 0.6868 & 0.2685 \\ 0.2685 & 0.3270 \end{pmatrix}, \quad R_E = \det = 0.1525\ldots
\]
The BSD formula predicts $L''(E,1)/2 = \Omega_E \cdot R_E \cdot \prod c_p \cdot |\Sha|/|E_{\mathrm{tors}}|^2$. With $\Omega_E = 4.9804\ldots$, $c_{389} = 1$, $|E_{\mathrm{tors}}| = 1$, and the reference value $|\Sha| = 1$, this gives $L''(E,1)/2 = 0.7593\ldots$, matching numerical computation. This is a rank-2 consistency check, not bridge evidence: Axioms~I and~IV are not rigorously proven in this range, and the input uses known BSD data.
\end{example}


% ===================================================================
% 6. HIGHER RANK
% ===================================================================
\section{Higher Rank: The Open Frontier}
\label{sec:higher_rank}

\subsection{The Height Saturation Conjecture}

For general rank $r$, we formulate:

\begin{conjecture}[Height Saturation -- Reformulation of BSD in positivity language]
\label{conj:height_saturation}
Let $E/\mathbb{Q}$ with $\rank E(\mathbb{Q}) = r$. Let $P_1, \ldots, P_r$ be a basis of $E(\mathbb{Q})/E(\mathbb{Q})_{\mathrm{tors}}$. Then:
\begin{equation}
L^{(r)}(E,1) \sim R_E = \det\!\left(\langle P_i, P_j \rangle\right) > 0\,,
\label{eq:height_saturation}
\end{equation}
in the sense that the ratio $L^{(r)}(E,1)/(r! \cdot \Omega_E \cdot R_E \cdot \prod c_p)$ equals $|\Sha|/|E_{\mathrm{tors}}|^2$.
\end{conjecture}

\begin{remark}[Relation to the classical BSD conjecture]
Conjecture~\ref{conj:height_saturation} is \textit{equivalent} to the strong BSD conjecture (Conjecture~\ref{conj:bsd_strong}): the ``height saturation'' formulation simply rephrases the BSD leading-term formula in the language of this paper's positivity framework. The name ``Height Saturation'' emphasizes the role of the positive-definite regulator $R_E > 0$ as the dominant geometric factor, but the mathematical content is identical to the classical formulation. We retain this reformulation because it highlights the positivity structure, not because it adds new mathematical content.
\end{remark}

This is the ``higher Gross-Zagier'' in its most optimistic form: the $r$-th derivative of $L$ is proportional to the $r \times r$ regulator -- the determinant of a positive-definite matrix.

\begin{remark}[Cauchy coefficient formulation]
\label{conj:height-saturation-identity}
The BSD leading-term formula can be expressed with the Cauchy coefficient formula, making the role of the Taylor coefficient explicit. This is a standard complex-analytic identity, not an independent conjecture. For an elliptic curve $E/\mathbb{Q}$ of analytic rank $r = \mathrm{ord}_{s=1} L(E,s)$:
\[
\begin{aligned}
\frac{L^{(r)}(E, 1)}{r!}
&= \frac{1}{2\pi i}
   \oint_{|s-1|=\varepsilon}
   \frac{L(E, s)}{(s-1)^{r+1}} \, ds \\
&= \Omega_E \cdot \det(\langle P_i, P_j \rangle_{\mathrm{NT}})
   \cdot \frac{\#\mathrm{Sha}(E) \cdot \prod c_p}
   {\#E(\mathbb{Q})_{\mathrm{tors}}^2}.
\end{aligned}
\]
where $P_1, \ldots, P_r$ are generators of $E(\mathbb{Q}) / E(\mathbb{Q})_{\mathrm{tors}}$, and $\varepsilon>0$ is chosen small enough that the circle contains no singularity of the completed local expression other than the expansion point. The unweighted integral $\oint L(E,s)\,ds$ would be zero, since $L(E,s)$ is holomorphic at $s=1$.
\end{remark}

\begin{remark}[Normalisation of the coefficient extraction]\label{rem:disc-normalisation}
Remark~\ref{conj:height-saturation-identity} uses the standard Cauchy normalisation: if $L(E,s)=\sum_{k\geq r} a_k(s-1)^k$ near $s=1$, then
\[
\frac{1}{2\pi i}\oint_{|s-1|=\varepsilon}
  \frac{L(E,s)}{(s-1)^{r+1}}\,ds
 = a_r = \frac{L^{(r)}(E,1)}{r!}.
\]
Equivalently, one may use the Fourier coefficient
$\varepsilon^{-r}(2\pi)^{-1}\int_0^{2\pi} L(E,1+\varepsilon e^{i\theta})e^{-ir\theta}\,d\theta$.
Under BSD, this coefficient equals $\Omega_E \cdot R_E \cdot \prod c_p \cdot |\mathrm{Sha}| / |E_{\mathrm{tors}}|^2$.
\end{remark}

\begin{remark}[Height saturation as a selection principle]
\label{rem:height-saturation}
The conjecture reformulates the BSD leading-term formula as a coefficient identity rather than a pointwise evaluation: the Cauchy functional extracts the leading Taylor coefficient, which BSD equates to a product of arithmetic invariants dominated by the positive-definite regulator $R_E$. This perspective suggests viewing the BSD formula as a \emph{selection principle}: among all leading-term structures \textit{of the form ``product of canonical arithmetic invariants''} (see the heuristic argument of Proposition~\ref{thm:uniqueness}), the positivity of the N\'eron-Tate height pairing makes the BSD formula the unique natural candidate, with $\Sha$ providing the remaining arithmetic correction. The rigorous formulation of this uniqueness is provided by the Bloch-Kato conjecture \cite{BlochKato1990}.
\end{remark}

\begin{remark}[Numerical diagnostics via LMFDB]\label{rem:lmfdb-saturation}
LMFDB and Cremona data allow consistency diagnostics, not a replacement for the missing algebraic control. For curves with available analytic leading terms and regulator data one can form
\[
Q_E =
\frac{L^{(r)}(E,1)\,\#E(\mathbb{Q})_{\mathrm{tors}}^2}
     {r!\,\Omega_E\,R_E\,\prod c_p}.
\]
BSD predicts that $Q_E=\#\mathrm{Sha}(E)$, hence a non-negative square order when $\Sha$ is finite. Such computations test the normalisation, local factors, regulator data and sample-level consistency with BSD. They do not establish (H1), (H2), or (H3), and they must be kept separate from any proof claim. When the input uses a known Mordell--Weil basis, known rank, or regulator data, the computation is reference-normalised and must be flagged as advice rather than bridge evidence. Smith's distribution results for $2^\infty$-Selmer groups \cite{Smith2022} provide a probabilistic background for designing rank-stratified samples.
\end{remark}

\begin{remark}[Rank-2 height saturation]
\label{rem:rank2-saturation}
For 17 rank-2 elliptic curves from the Cremona database (conductors $389 \leq N \leq 5077$), the regulator satisfies $R > 0$ in all cases ($R_{\min} = 0.095$, $R_{\max} = 1.700$). Log-log regression yields $R \sim N^{0.099}$, consistent with sublinear growth. \textbf{Caveat:} The sample size ($n = 17$) is too small for statistically significant conclusions; these results should be viewed as preliminary consistency checks, not as statistical confirmation of a scaling law. The positivity $R > 0$ is guaranteed a priori by the positive-definiteness of the N\'eron-Tate height (Theorem~\ref{thm:NT_positive}). Script: \texttt{compute\_rank2\_lmfdb.py}.
\end{remark}

\subsection{What would a proof require?}

A proof of Conjecture~\ref{conj:height_saturation} via the positivity framework requires three hypotheses, which we label for reference throughout the paper:

\textbf{(H1) Higher Gross-Zagier formula.} An explicit identity:
\begin{equation}
L^{(r)}(E,1) = C_r \cdot \det\!\left(\langle P_i, P_j \rangle\right)\,,
\label{eq:higher_gz}
\end{equation}
where $P_1, \ldots, P_r$ are canonical arithmetic points (``higher Heegner points'') and $C_r > 0$ is an explicit constant.

\textit{Status:} For $r = 1$, this is Gross-Zagier. For $r = 2$, partial results exist via diagonal restriction of $p$-adic families \cite{DarmonRotger2017}. Kim \cite{Kim2024} has obtained a higher Gross-Zagier formula refining the connection between Heegner point Kolyvagin systems and Selmer group structure. For general $r$, this is a major open problem. Recent work on Euler systems for higher-rank motives \cite{LSZ2020} provides progress but not a complete formula.

\textbf{(H2) Higher Kolyvagin system.} An Euler system machinery that bounds $\Sel(E/\mathbb{Q})$ using $r$ classes simultaneously, establishing:
\begin{equation}
\rank E(\mathbb{Q}) \leq \ord_{s=1} L(E,s)\,.
\end{equation}

\textit{Status:} For $r = 1$, this is Kolyvagin's method. For general $r$, Mazur-Rubin \cite{MazurRubin2004} have developed the theory of Kolyvagin systems abstractly, and Burns-Sano \cite{BurnsSano2021} and Burns-Sakamoto-Sano \cite{BurnsSakamotoSano2025} have extended this to higher rank Euler, Kolyvagin, and Stark systems over general coefficient rings. Kim \cite{Kim2024} has obtained a Kolyvagin system-theoretic refinement of the Gross-Zagier formula. However, the \textit{existence} of a Kolyvagin system of sufficient depth for arbitrary rank is not established.

\textbf{(H3) $\Sha$ finiteness for general rank.} Axiom~I must hold unconditionally:
\begin{equation}
|\Sha(E/\mathbb{Q})| < \infty \quad \forall\, E/\mathbb{Q}\,.
\end{equation}

\textit{Status:} Open for $\rank \geq 2$. This is widely expected but no general proof exists.

\subsection{The positivity gap}

The situation can be summarized:

\begin{center}
\begin{tabular}{lccc}
\toprule
Axiom & Rank 0 & Rank 1 & Rank $\geq 2$ \\
\midrule
I (Finite $\Sha$) & $\checkmark$ & $\checkmark$ & Open \\
II (Height pos.) & $\checkmark$ & $\checkmark$ & $\checkmark$ \\
III (Iwasawa) & $\checkmark$ & $\checkmark$ & Partial \\
IV (Euler sys.) & $\checkmark$ & $\checkmark$ & Open \\
\bottomrule
\end{tabular}
\end{center}

The ``positivity'' (Axiom~II) is always available -- it is a theorem. The obstruction lies in the ``composition'' (Axioms~I, III, IV): the Euler-system/Kolyvagin/Iwasawa machinery is not yet fully developed for higher rank.

In summary: the positivity (Axiom~II) is established, but the composition law (Axioms~I, III, IV) is incomplete for rank~$\geq 2$. The BSD formula is forced once the Euler-system composition is established for arbitrary rank. This paper is part of a broader programme \cite{Geiger2026e}; the connections are purely motivational and not used in any proof.


% ===================================================================
% 7. IWASAWA THEORY AND COMPOSITION
% ===================================================================
\section{Iwasawa Theory as Composition Coherence}
\label{sec:iwasawa}

\subsection{\texorpdfstring{The $p$-adic perspective}{The p-adic perspective}}

Iwasawa theory provides the clearest realization of Axiom~III (controlled growth) and connects it naturally to Axiom~IV (Euler-system coherence). The main conjecture for elliptic curves \cite{Kato2004,SkinnerUrban2014} is established in the following specific cases:
\begin{itemize}
\item \textbf{Kato's divisibility (one direction):} For all primes $p$ where $E$ has good ordinary reduction and $\bar{\rho}_{E,p}$ is surjective. Applies to \textit{all ranks}, but only gives the inclusion $\mathrm{char}_\Lambda(\Sel^\vee) \subseteq (\mathcal{L}_p)$, i.e., the characteristic ideal of the Selmer dual is contained in the ideal generated by the $p$-adic $L$-function, providing an upper bound on the Selmer module.
\item \textbf{Skinner-Urban (full equality):} For ordinary primes $p$ with $p \nmid a_p$, provided $\bar{\rho}_{E,p}$ is irreducible and additional local hypotheses hold. Does \textit{not} cover supersingular primes.
\item \textbf{Supersingular primes:} Require the signed Selmer group formulation (Kobayashi, Pollack, Sprung); the main conjecture in this setting remains largely open.
\item \textbf{Not covered:} Primes of bad additive reduction, primes where $\bar{\rho}_{E,p}$ is reducible, and the general case without residual irreducibility.
\end{itemize}
The main conjecture asserts:
\begin{equation}
\mathrm{char}_{\Lambda}\!\left(\Sel_{p^\infty}(E/\mathbb{Q}_\infty)^\vee\right) = (L_p(E))\,,
\label{eq:iwasawa_mc}
\end{equation}
where $\Lambda = \mathbb{Z}_p[[T]]$ is the Iwasawa algebra, $\Sel_{p^\infty}(E/\mathbb{Q}_\infty)^\vee$ is the Pontryagin dual of the Selmer group over the cyclotomic $\mathbb{Z}_p$-extension, and $L_p(E)$ is the $p$-adic $L$-function.

\begin{proposition}[Iwasawa main conjecture as composition law]
\label{prop:iwasawa_composition}
The Iwasawa main conjecture provides the composition structure in the following precise sense: let $\mathrm{char}_\Lambda(M)$ denote the characteristic ideal of a finitely generated torsion $\Lambda$-module $M$. The main conjecture identifies $\mathrm{char}_\Lambda(\Sel_{p^\infty}(E/\mathbb{Q}_\infty)^\vee) = (L_p(E))$. The ``composition'' is the control theorem (Mazur): specialising at layer $n$ of the $\mathbb{Z}_p$-tower, the Selmer group at level $n$ is related to the Selmer group at level $n-1$ by an exact sequence whose kernel and cokernel are controlled by the Euler factor at $p$. This layer-by-layer compatibility is the arithmetic analogue of norm-compatibility in Euler systems.
\end{proposition}

The algebraic structure is: $\Lambda = \mathbb{Z}_p[[T]]$ is a power series ring, and the characteristic ideal is generated by a single element $f_E(T)$. The specialisation maps $\chi_n: \Lambda \to \mathbb{Z}_p[\Gal(\mathbb{Q}_n/\mathbb{Q})]$, obtained by evaluating at $T = \zeta_{p^n} - 1$, form a projective system:
\begin{equation}
\chi_m = \chi_n \circ \mathrm{pr}_{m,n} \quad \text{for } m \geq n\,,
\end{equation}
where $\mathrm{pr}_{m,n}$ is the natural projection. The control theorem ensures that $\Sel_{p^\infty}(E/\mathbb{Q}_n)$ is determined by $\Sel_{p^\infty}(E/\mathbb{Q}_\infty)$ via restriction with bounded kernel and cokernel. This controlled compatible system is the $p$-adic composition law constraining the arithmetic of $E$ across the tower.

\subsection{\texorpdfstring{The $\mu = 0$ conjecture}{The mu = 0 conjecture}}

The Iwasawa $\mu$-invariant measures the ``$p$-adic mass'' of $\Sha$ in the tower. The conjecture $\mu = 0$ (proven by Ferrero-Washington for abelian fields, expected in general) ensures that no infinite $p$-power contribution hides in $\Sha$. This is precisely Axiom~I restricted to the $p$-adic tower.

\subsection{\texorpdfstring{Specialization to $s = 1$}{Specialization to s = 1}}

The key connection between Iwasawa theory and BSD is the \textit{specialization}: evaluating the $p$-adic $L$-function at $T = 0$ (or its derivatives) recovers the complex $L$-function values:
\begin{equation}
\begin{aligned}
L_p(E, 0) &\sim L(E, 1)/\Omega_E^+ \\
&\quad (\text{up to periods and Euler factors})\,.
\end{aligned}
\end{equation}

For rank $r$, the relevant specialization involves the $r$-th derivative of $L_p(E, T)$ at $T = 0$. The Iwasawa main conjecture then relates this to the characteristic ideal of the Selmer module, providing a $p$-adic version of the BSD formula.

\subsection{\texorpdfstring{The $p$-adic to archimedean bridge}{The p-adic to archimedean bridge}}
\label{sec:p-adic-bridge}

The Iwasawa main conjecture and Euler-system machinery (Axioms~III--IV) live naturally in the $p$-adic world, while the BSD conjecture concerns the \textit{complex} $L$-function $L(E,s)$ at the archimedean place $s = 1$. The passage between these two domains is controlled by Fontaine's comparison isomorphism $H^1_{\mathrm{dR}}(E/\mathbb{Q}_p) \cong H^1_{\mathrm{\acute{e}t}}(E, \mathbb{Q}_p) \otimes B_{\mathrm{dR}}$ from $p$-adic Hodge theory, which provides a dictionary between $p$-adic periods and complex periods ($\Omega_E$), between $p$-adic heights ($\hat{h}_p$) and archimedean heights ($\hat{h}$), and between $p$-adic $L$-values ($L_p(E,0)$) and complex $L$-values ($L(E,1)$).

For rank~0, the bridge is well understood: at a prime $p$ of good ordinary reduction, the interpolation formula gives $L_p(E,0) = (1 - \alpha_p^{-1})^2 \cdot L(E,1)/\Omega_E^+$ up to a $p$-adic unit, where $\alpha_p$ is the unit root of $X^2 - a_p X + p$ (for multiplicative reduction, the Mazur-Tate-Teitelbaum conjecture, proven by Greenberg-Stevens \cite{GreenbergStevens1993}, introduces the $\mathcal{L}$-invariant). For rank~1, the Gross-Zagier and Perrin-Riou \cite{PerrinRiou1987} formulas provide compatible archimedean and $p$-adic height couplings, closing the bridge at the level of first derivatives.

For rank $\geq 2$, the bridge is the deepest obstruction in the positivity framework. Even if a $p$-adic height coupling $L_p^{(r)}(E,0) \sim \det(\hat{h}_p(P_i, P_j))$ were established via diagonal cycles \cite{DarmonRotger2017}, the translation to the archimedean formula $L^{(r)}(E,1) \sim R_E$ requires two additional ingredients: (i) a comparison between $p$-adic and archimedean regulators, $R_{E,p} \sim R_E$, which is known up to a $p$-adic unit for CM curves but open in general; and (ii) an interpolation formula for the $r$-th derivative generalizing Mazur-Tate-Teitelbaum beyond rank~0.

This bridge crystallizes the structural tension in our framework: the \textit{positivity} (Axiom~II) is an archimedean phenomenon -- $\hat{h}$ is a real-valued positive-definite form -- while the \textit{composition} (Axioms~III--IV) is a $p$-adic phenomenon -- Euler systems and Iwasawa modules live over $\mathbb{Z}_p$. Unifying these two ``positivity domains'' via the comparison isomorphism is the key technical challenge for extending the BSD positivity program beyond rank~1.


% ===================================================================
% 8. THE PRECISE OBSTRUCTION
% ===================================================================
\section{The Precise Obstruction}
\label{sec:obstruction}

\subsection{What blocks a full proof}

The positivity framework identifies a precise bottleneck: the \textit{absence of a universal higher Gross-Zagier formula}. Specifically:

\begin{enumerate}
\item \textbf{For rank 1:} The Heegner point $P_K$ provides a canonical ``geometric direction'' whose height couples to $L'(E,1)$. The composition law (Kolyvagin's Euler system) then excludes alternatives.

\item \textbf{For rank $r$:} One would need $r$ independent canonical points $P_1, \ldots, P_r$ whose height \textit{matrix} couples to $L^{(r)}(E,1)$. The candidates are:
\begin{itemize}
\item Higher Heegner points on Shimura varieties of dimension $> 1$.
\item Generalized Heegner cycles (Bertolini-Darmon-Prasanna \cite{BDP2013}).
\item Diagonal restriction of $p$-adic families (Darmon-Rotger \cite{DarmonRotger2017}).
\end{itemize}
None currently provides a complete ``higher Gross-Zagier formula'' in the form~\eqref{eq:higher_gz}.
\end{enumerate}

\subsection{The obstruction in positivity language}

The obstruction can be stated precisely:

\begin{quote}
\textit{The positivity of $R_E = \det(\langle P_i, P_j \rangle) > 0$ is a theorem (from $\hat{h}$ being positive-definite). The missing piece is a formula equating $L^{(r)}(E,1)$ to $R_E$ times computable constants -- i.e., the explicit coupling of the analytic and geometric positive structures.}
\end{quote}

Once such a coupling is established, the No-Go argument closes:
\begin{enumerate}
\item $L^{(r)}(E,1) = C_r \cdot R_E > 0$ (positivity of heights), hence $\ord_{s=1} L(E,s) \leq r$.
\item The Euler-system bound (Axiom~IV) gives $\rank E(\mathbb{Q}) \leq \ord_{s=1} L(E,s)$, hence $\ord_{s=1} L(E,s) \geq r$.
\item Together: $\ord_{s=1} L(E,s) = r = \rank E(\mathbb{Q})$.
\end{enumerate}

The entire argument rests on coupling two positive structures -- one analytic ($L$-derivatives), one geometric (height pairing) -- and showing they must ``saturate'' at the same order.


% ===================================================================
% 9. CONNECTION TO THE BROADER PROGRAM
% ===================================================================
\section{Connection to the Broader Program}
\label{sec:connection}

\subsection{The positivity pattern across mathematics}

The BSD framework fits into the same positivity pattern observed in:

\begin{center}
\footnotesize
\resizebox{\columnwidth}{!}{%
\begin{tabular}{@{}llll@{}}
\toprule
Problem & Positive form & Composition & Status \\
\midrule
Weil conj. & Frobenius ev. & \'Etale coh. & Proven \\
Hodge conj. & HR bilinear & Galois-HR & Open \\
\textbf{BSD} & \textbf{N\'eron-Tate} & \textbf{Euler sys.} & \textbf{Rank $\leq 1$} \\
\bottomrule
\end{tabular}
}
\end{center}

In each case:
\begin{itemize}
\item The \textit{positivity} is the ``easy'' part (a theorem or a well-motivated condition).
\item The \textit{composition/coherence} is the ``hard'' part (requires deep structural work).
\item The \textit{conjunction} of positivity and composition forces the desired result.
\end{itemize}

\subsection{Link to the Hodge approach}

The N\'eron-Tate height $\hat{h}$ on $E(\mathbb{Q})$ and the Hodge-Riemann form $Q_L$ on $H^{p,p}(X)$ play structurally identical roles:
\begin{itemize}
\item Both are canonical positive-definite bilinear forms.
\item Both detect ``arithmetic directions'' (rational points / algebraic cycles).
\item Both require a ``composition law'' (Euler systems / Galois-HR stability) to force the desired conclusion.
\end{itemize}

The Hodge analog of the Gross-Zagier formula would be: an explicit identity coupling the ``analytic'' Hodge-Riemann form to the ``geometric'' cycle class map, showing that positivity of one forces algebraicity of the other.

\subsection{Link to the Szpiro conjecture}
\label{sec:szpiro}

The BSD framework connects to the Szpiro conjecture (equivalent to the $abc$ conjecture) through the Tamagawa numbers $c_p$ in the leading-term formula~\eqref{eq:bsd_formula}.

\subsubsection{Szpiro as a local discriminant-exponent bound}

\begin{conjecture}[Szpiro {\cite{Szpiro1981}}]
\label{conj:szpiro}
For every $\varepsilon > 0$, there exists $C_\varepsilon > 0$ such that for every elliptic curve $E/\mathbb{Q}$ with minimal discriminant $\Delta_{\min}$ and conductor $N$:
\begin{equation}
\log|\Delta_{\min}| \leq (6 + \varepsilon) \log N + C_\varepsilon\,.
\label{eq:szpiro}
\end{equation}
\end{conjecture}

For a semistable curve $E/\mathbb{Q}$, all bad reduction is multiplicative.  Write
\[
  n_p := v_p(\Delta_{\min}) .
\]
Then $N = \prod_{p \mid \Delta} p$ and
\begin{equation}
\sum_{p \mid N} n_p \log p
= \log|\Delta_{\min}|
\leq (6+\varepsilon)\sum_{p \mid N}\log p + C_\varepsilon\,.
\label{eq:szpiro_tamagawa}
\end{equation}
For split multiplicative reduction at~$p$, the Tamagawa number equals this local exponent, $c_p=n_p$ (\cite{Silverman2009}, Table~IV.9.1).  For nonsplit multiplicative reduction, however, $c_p$ may be only $1$ or $2$ while $n_p$ is larger.  Thus Szpiro is fundamentally a \emph{prime-weighted local discriminant-exponent} bound; only in the split multiplicative subcase can it be read literally as a prime-weighted Tamagawa-exponent bound.  This is stronger and more structured than an unweighted bound on $\prod_p c_p$ (or even on $\prod_p n_p$): a large exponent at a large prime carries more discriminant mass than the same exponent at a small prime.

\subsubsection{\texorpdfstring{Diagnostic chain: BSD $+$ bounds expose the missing weighted step}{Diagnostic chain: BSD + bounds expose the missing weighted step}}

The BSD formula~\eqref{eq:bsd_formula} yields:
\begin{equation}
\prod_p c_p = \frac{L^{(r)}(E,1)}{r!} \cdot \frac{|E_{\mathrm{tors}}|^2}{\Omega_E \cdot R_E \cdot |\Sha|}\,.
\label{eq:tamagawa_from_bsd}
\end{equation}
This identity controls the \emph{unweighted product} of Tamagawa factors.  It does not by itself imply Szpiro, because Szpiro controls the prime-weighted sum in~\eqref{eq:szpiro_tamagawa}.  A route from BSD to Szpiro therefore requires the following ingredients plus an additional weighted-local upgrade:
\begin{enumerate}
\item[\textbf{(B1)}] \textbf{Upper bound on $L$-values:} $L^{(r)}(E,1)/r! \leq C \cdot N^{r/2 + \varepsilon}$.  Unconditional for $r = 0$ (convexity bound); conditional on GRH for general~$r$.
\item[\textbf{(B2)}] \textbf{Lower bound on the period:} $\Omega_E \geq c \cdot N^{-1/2 - \varepsilon}$.  Follows from bounds on the modular parametrisation degree.
\item[\textbf{(B3)}] \textbf{Lower bound on the regulator:} $R_E \geq c \cdot (\log N)^r$.  Requires an effective lower bound on the N\'eron-Tate height of non-torsion points.
\item[\textbf{(B4)}] \textbf{Torsion bound:} $|E_{\mathrm{tors}}|^2 \leq 256$.  This is Mazur's torsion theorem ($|E(\mathbb{Q})_{\mathrm{tors}}| \leq 16$).
\item[\textbf{(B5)}] \textbf{Weighted local upgrade:} a mechanism converting unweighted control of $\prod_p c_p$ into control of the actual local exponents $\sum_{p\mid N}n_p\log p$, including the split/nonsplit multiplicative gap, or directly bounding each local discriminant contribution in terms of conductor data.
\end{enumerate}
Combining (B1)--(B4) under the BSD assumption yields only:
\begin{equation}
\prod_p c_p \leq C' \cdot \frac{N^{r/2 + 1/2 + 2\varepsilon}}{(\log N)^r}\,.
\label{eq:tamagawa_bound}
\end{equation}
For semistable curves of rank~$0$, this gives a strong unweighted Tamagawa-product constraint.  It is \emph{not} yet a Szpiro bound unless (B5) supplies both the missing prime weights and the local passage from Tamagawa factors to the actual discriminant exponents $n_p$.

\begin{remark}[Circularity in step (B3)]
\label{rem:szpiro_circularity}
Even before the weighted upgrade, the chain has a second obstruction: circularity in (B3).  The best known effective lower height bound is Silverman's $\hat{h}(P) \geq c \cdot \log N$ for non-torsion $P$~\cite{Silverman2009}, which is \textit{conditional on the $abc$ conjecture}.  Since $abc \Leftrightarrow$ Szpiro, the logical structure would be:
\[
\begin{gathered}
\text{BSD} + \underbrace{\text{abc}}_{\text{used in (B3)}} \\
+ \underbrace{\text{weighted local upgrade}}_{\text{missing (B5)}}
\Longrightarrow
\underbrace{\text{Szpiro}}_{\Leftrightarrow\, \text{abc}}\,.
\end{gathered}
\]
Breaking this circularity requires a height bound that does not depend on the discriminant or on $abc$, and a separate mechanism for the prime-weighted local exponents.
\end{remark}

\subsubsection{A theta-product bridge}

The following identity, connecting theta values at torsion points to the Dedekind eta function (and hence to the modular discriminant), provides a potential non-circular link between torsion structure and the discriminant.

% TOOL-NOTE: The theta-eta product identity below is consolidated in
%   FST_MATHEMATICS/_tools/THETA_ETA_IDENTITY.md
% Identical statement and binding notation. The same identity is used
% in abc/abc_Theta_Tamagawa_EN.tex Proposition 4.1.

\begin{proposition}[Theta-eta product identity at $l$-torsion]
\label{prop:theta_eta}
For any odd prime $l$ and $\tau \in \mathcal{H}$, with $q = e^{2\pi i \tau}$:
\begin{equation}
\prod_{j=1}^{l-1} \theta_1\!\left(\frac{j}{l} \,\bigg|\, \tau\right) = l \cdot \eta(\tau)^{l-3} \cdot \eta(l\tau)^2\,,
\label{eq:theta_eta_identity}
\end{equation}
where $\theta_1(z|\tau) = -i \sum_{n \in \mathbb{Z}} (-1)^n q^{(n+1/2)^2/2} e^{(2n+1)\pi i z}$ is the Jacobi theta function and $\eta(\tau) = q^{1/24} \prod_{n=1}^{\infty}(1-q^n)$ is the Dedekind eta function.
\end{proposition}

\begin{proof}
The Jacobi triple product gives $\theta_1(z|\tau) = -2q^{1/8}\sin(\pi z)\prod_{n=1}^{\infty}(1-q^n)(1-e^{2\pi i z}q^n)(1-e^{-2\pi i z}q^n)$.
Setting $\omega = e^{2\pi i/l}$, the cyclotomic identity $\prod_{j=1}^{l-1}(1-\omega^j x) = 1 + x + \cdots + x^{l-1}$ applied to each factor of the infinite product, together with $\prod_{j=1}^{l-1}\sin(\pi j/l) = l/2^{l-1}$ (the Chebyshev sine identity), yields:
\[
\prod_{j=1}^{l-1}\theta_1(j/l|\tau)
= l\,q^{(l-1)/8}
\prod_{n=1}^{\infty}(1-q^n)^{l-3}(1-q^{ln})^2\,.
\]
Since $\eta(\tau) = q^{1/24}\prod_{n\geq1}(1-q^n)$ and $\eta(l\tau) = q^{l/24}\prod_{n\geq1}(1-q^{ln})$, the eta-product $l\eta(\tau)^{l-3}\eta(l\tau)^2$ has exactly the prefactor $lq^{(l-3)/24+2l/24}=lq^{(l-1)/8}$ and the same infinite product. This proves the identity.
\end{proof}

\begin{remark}[Connection to the discriminant]
\label{rem:theta_discriminant}
Since $\eta(\tau)^{24} = \Delta(\tau)/(2\pi)^{12}$ (where $\Delta$ is the modular discriminant), identity~\eqref{eq:theta_eta_identity} gives:
\begin{equation}
\begin{aligned}
\sum_{j=1}^{l-1} \log|\theta_1(j/l|\tau)|
&= \log l
 + \frac{l-3}{24}\log\frac{|\Delta(\tau)|}{(2\pi)^{12}} \\
&\quad + \frac{2}{24}\log\frac{|\Delta(l\tau)|}{(2\pi)^{12}}\,.
\end{aligned}
\label{eq:theta_log_discriminant}
\end{equation}
For large $l$, the dominant term is $(l/24)\log|\Delta(\tau)|$: the product of theta values at $l$-torsion points encodes the discriminant.
\end{remark}

\begin{remark}[Tate curve application and the Szpiro connection]
\label{rem:tate_szpiro}
For an elliptic curve $E/\mathbb{Q}_p$ with split multiplicative reduction (Tate curve $E_q \cong \mathbb{G}_m/q^{\mathbb{Z}}$), the $l$-torsion points correspond to $\zeta_l^j \cdot q^{k/l}$ for $0 \leq j,k < l$.  The identity~\eqref{eq:theta_eta_identity} then relates the $l$-torsion structure of $E$ to the local discriminant exponent $n_p=v_p(\Delta_{\min})$ (equivalently $c_p$ at split multiplicative places) \textit{without invoking height bounds or the $abc$ conjecture}.  This suggests a potential non-circular path to Szpiro: if the $l$-torsion theta values can be bounded in terms of the conductor~$N$ -- for instance via the Galois representation $\bar{\rho}_{E,l}$ or the Weil explicit formula for $L(E,s)$ -- the resulting bound on $n_p$ would yield the local input needed for the Szpiro conjecture directly.

The missing step is a prime-weighted local conductor estimate of the form:
\begin{equation}
\sum_{p\mid N}\sum_{j=1}^{l-1}-\log|\theta_1(j/l|\tau_p)|_p
\leq C(l,\varepsilon)\cdot\log N\,,
\label{eq:theta_conductor_bound}
\end{equation}
which is not currently established but is the correct weighted target.  Bounding the theta values via the conductor connects naturally to the Galois action on $E[l]$, where Serre's modularity results and level-lowering provide the link between the $l$-torsion representation and the conductor.
\end{remark}

\begin{remark}[Relation to Mochizuki's IUT approach]
\label{rem:iut_comparison}
The theta-eta product identity~\eqref{eq:theta_eta_identity} operates on the same objects -- theta values at torsion points of Tate curves -- that appear in Mochizuki's Inter-Universal Teichm\"uller Theory (IUT).  However, the IUT approach uses the theta values $q^{j^2}$ (quadratic in $j$) arising from the multiradial representation, whereas the identity~\eqref{eq:theta_eta_identity} concerns $\theta_1(j/l|\tau)$ (linear in $j$).  This distinction is analysed in~\cite{GeigerIUT2026}, where it is shown that the $O(l)$-vs-$O(l^2)$ scaling mismatch between the linear and quadratic theta values constitutes a structural obstruction to reproducing the IUT bound via classical methods.  The present approach differs fundamentally: it does not attempt to reproduce the IUT mechanism, but instead uses the product identity as a direct bridge from torsion to discriminant within the BSD framework.
\end{remark}


% ===================================================================
% 10. DISCUSSION
% ===================================================================
\section{Discussion}
\label{sec:discussion}

\subsection{What has been achieved}

\begin{enumerate}
\item \textbf{Axiomatization.} Four axioms -- finite capacity, height positivity, controlled Iwasawa growth, Euler-system composition -- that conditionally identify the BSD formula as the expected product normal form once (H1)--(H3) are supplied.

\item \textbf{Rank $\leq 1$ as a completed instance.} The Gross-Zagier + Kolyvagin argument is naturally a positivity No-Go proof: it couples an analytic derivative to a non-negative height and uses composition (Euler systems) to exclude alternatives.

\item \textbf{Higher rank identified.} The obstruction to higher rank is precisely the absence of (H1) a universal higher Gross-Zagier formula and (H2) a complete Kolyvagin system -- i.e., the composition law is incomplete.

\item \textbf{Exclusion of alternatives.} No alternative leading-term formula is compatible with all four axioms, \textit{provided} a coupling between $L^{(r)}(E,1)$ and the height pairing exists.
\end{enumerate}

\subsection{Limitations}

\begin{enumerate}
\item \textbf{The paper does not prove BSD.} It reformulates it as a positivity normal-form statement and identifies the missing bridge modules.

\item \textbf{Axiom~I ($\Sha$ finiteness) is conjectural for rank $\geq 2$.} This is itself a major open problem, though widely expected.

\item \textbf{The ``higher Gross-Zagier'' is speculative.} No complete formula of the form~\eqref{eq:higher_gz} exists for $r \geq 2$. Recent progress (Darmon-Rotger, Loeffler-Skinner-Zerbes) is promising but incomplete.

\item \textbf{The axiomatic framework is descriptive, not constructive.} The axioms identify what must be true, not how to prove it.
\end{enumerate}

\subsection{Outlook}

\begin{enumerate}
\item \textbf{$p$-adic BSD via Iwasawa theory.} The strongest near-term progress likely comes from extending the Iwasawa main conjecture to higher-rank settings, using the Euler systems of Loeffler-Skinner-Zerbes \cite{LSZ2020} and the higher rank Kolyvagin-Stark system theory of Burns-Sakamoto-Sano \cite{BurnsSano2021,BurnsSakamotoSano2025}.

\item \textbf{Higher Gross-Zagier.} Darmon's program of ``Stark-Heegner points'' and diagonal cycles provides candidate formulas for rank 2. Establishing these rigorously would close the positivity loop for $r = 2$.

\item \textbf{Computational verification.} The BSD formula has been verified numerically for thousands of curves of rank $\leq 3$ (see \cite{LMFDB}). Extending these computations to higher precision and higher rank tests the framework's predictions.

\item \textbf{Statistical BSD.} The work of Bhargava-Shankar \cite{BhargavaShankar2015} shows that a positive proportion of elliptic curves over $\mathbb{Q}$ have rank~0 and satisfy the BSD conjecture. These statistical results provide further evidence that the positivity framework captures the generic behavior.
\end{enumerate}


\begin{remark}[Arakelov bridge to Hodge]
\label{rem:arakelov-bridge}
The N\'eron--Tate height pairing encodes positivity at finite (non-archimedean) primes. Together with the Hodge--Riemann form (positivity at the infinite prime, cf.\ the companion Hodge paper), the Arakelov intersection pairing provides the natural unified framework in which both are special cases of a single positivity principle.
\end{remark}


\begin{remark}[Geometric interpretation of the regulator]
\label{rem:entropic-capacity}
The regulator $R_E = \det(\langle P_i, P_j \rangle_{\mathrm{NT}})$ is the covolume of the Mordell--Weil lattice $E(\mathbb{Q})/E(\mathbb{Q})_{\mathrm{tors}}$ in the N\'eron--Tate metric. The BSD formula equates the analytic datum $L^{(r)}(E,1)$ to this geometric volume (times local correction factors). A connection to random matrix theory exists via the Keating--Snaith conjectures on the distribution of $L$-function derivatives, but we do not develop this here.
\end{remark}

\subsection{Research directions for higher rank}
\label{sec:higher_rank_directions}

The following remarks survey concrete approaches toward establishing the missing hypotheses (H1)--(H3) for rank~$\geq 2$.

\subsubsection{Diagonal cycles and higher Gross-Zagier formulas}

\begin{remark}[Diagonal-cycle / Euler-system hybrid for rank-2 CM curves]
\label{rem:diagonal-euler}
For an elliptic curve $E/\mathbb{Q}$ with CM by an imaginary quadratic field $K$, the Selmer group $\Sel_{p^\infty}(E/K)$ decomposes into $\pm$-eigenspaces under complex conjugation. If $\rank E(\mathbb{Q}) = 2$ and $E$ has CM, one could attempt to construct two independent cohomology classes -- one from a Heegner-type cycle and one from a Beilinson--Flach element -- whose heights span the full regulator matrix. The key challenge is establishing the non-degeneracy of the resulting $2 \times 2$ height matrix: $\det(\langle c_i, c_j \rangle) > 0$ must be shown to couple to $L''(E,1)$. Darmon and Rotger~\cite{DarmonRotger2017} provide partial results in this direction via diagonal cycles on triple products of modular curves.
\end{remark}

\begin{remark}[Heegner-simplex approach for rank $r$]
\label{rem:heegner-simplex}
For rank~$r$, one seeks $r$-dimensional cycles on Kuga-Sato varieties or higher-dimensional Shimura varieties whose Abel-Jacobi images generate a full-rank subgroup of the relevant Selmer group. Concretely, one would need:
\begin{enumerate}
\item A Shimura variety $\mathrm{Sh}(G, X)$ of dimension $\geq r$ with special cycles $Z_1, \ldots, Z_r$ indexed by CM data;
\item Non-trivial Abel-Jacobi images $\mathrm{AJ}(Z_i) \in H^1_f(\mathbb{Q}, V_p(E))$;
\item An explicit formula: $L^{(r)}(E,1) = C_r \cdot \det(\langle \mathrm{AJ}(Z_i), \mathrm{AJ}(Z_j) \rangle_p)$.
\end{enumerate}
The group-theoretic framework (replacing $\mathrm{GL}_2$ by $\mathrm{GSp}_{2r}$ or a product of copies of $\mathrm{GL}_2$) is provided by Loeffler--Skinner--Zerbes~\cite{LSZ2020}, but the explicit height formula remains open.
\end{remark}

\begin{remark}[Plectic Gross-Zagier formula]
\label{rem:plectic-gz}
Nekov\'a\v{r} and Scholl developed a ``plectic'' extension of Hodge theory, introducing plectic Galois actions and plectic Hodge structures that generalize the classical Hodge decomposition to higher-rank settings. A \textit{plectic Gross-Zagier formula} would express $L^{(r)}(E,1)$ in terms of the plectic regulator -- a higher determinant constructed from the plectic Hodge structure on the motive of $E$. This program, still in its early stages, has been continued by Li-Huerta and others; a complete formula would directly establish hypothesis (H1) and close the positivity loop for general rank.
\end{remark}

\begin{remark}[$p$-adic deformation of non-diagonal Heegner cycles]
\label{rem:p-adic-deformation}
Darmon's program of Stark-Heegner points constructs $p$-adic points on $E$ from $p$-adic integration on Shimura curves, without relying on the Heegner hypothesis. For rank~$\geq 2$, one could attempt to $p$-adically deform families of non-diagonal cycles on higher-dimensional Shimura varieties (e.g., unitary Shimura varieties associated to $\mathrm{U}(2,1)$ or $\mathrm{U}(3)$), obtaining families of cohomology classes whose specializations provide the required $r$ independent elements. The key obstruction is rationality: the $p$-adic constructions do not automatically yield $\mathbb{Q}$-rational points, and a descent argument is needed.
\end{remark}

\subsubsection{Euler systems for higher rank}

\begin{remark}[Beilinson-Flach elements and Rankin-Selberg convolutions]
\label{rem:beilinson-flach}
Kings, Loeffler, and Zerbes~\cite{KLZ2017} construct Euler systems from Beilinson--Flach elements for Rankin-Selberg convolutions $L(f \otimes g, s)$, where $f$ and $g$ are modular forms. When specialized to the case $f = g$ (the symmetric square), these elements provide information about the $p^\infty$-Selmer group of $E$ beyond what Kato's Euler system yields. For rank-2 curves, one could attempt to construct a ``rank-2 Kolyvagin system'' from the combination of Kato's classes and Beilinson-Flach elements, each providing control over one eigenspace of the Selmer group. The $p$-adic Eisenstein elements in multi-dimensional Hida families~\cite{LLZ2014} provide the interpolation framework.
\end{remark}

\begin{remark}[Higher Kolyvagin systems and Selmer splitting]
\label{rem:kolyvagin-higher}
Mazur and Rubin~\cite{MazurRubin2004} developed the abstract theory of Kolyvagin systems: a module over the Iwasawa algebra satisfying certain local conditions that yield upper bounds on Selmer groups. Burns and Sano~\cite{BurnsSano2021} extended this to higher rank systems over general coefficient rings, and Burns, Sakamoto, and Sano~\cite{BurnsSakamotoSano2025} proved the existence of a canonical higher Kolyvagin derivative map between Euler and Kolyvagin systems in any given rank, as conjectured by Mazur and Rubin. For rank~$r$, one needs a ``rank-$r$ Kolyvagin system'' -- a collection of $r$ classes satisfying simultaneous norm-compatibility relations. For CM curves, the endomorphism ring provides a natural splitting of the Selmer group into eigenspaces (under the CM action), and one can attempt to construct independent Kolyvagin classes in each eigenspace. Numerical testing of ``rank-$r$ Kolyvagin relations'' -- verifying the predicted Selmer bounds for specific rank-2 and rank-3 curves with CM -- would provide evidence for this approach.
\end{remark}

\begin{remark}[Recent progress on rank $\geq 2$]\label{rem:recent-bsd}
Significant progress toward higher-rank BSD has been achieved in recent years. Burungale, Castella, and Kim~\cite{BurungaleCastellaKim2021} proved Perrin-Riou's Heegner point main conjecture, building on Howard's bipartite Euler systems and Wei Zhang's work on Kolyvagin's conjecture, establishing the Iwasawa main conjecture for the Tate-Shafarevich group over anticyclotomic extensions. Kim~\cite{Kim2024} obtained a higher Gross-Zagier formula by comparing Heegner point Kolyvagin systems with Kurihara numbers, refining the structural understanding of Selmer groups. In 2024, Loeffler and Zerbes \cite{LoefflerZerbes2024} constructed a universal Euler system for $\mathrm{GSp}(4)$, showing that all Euler system classes lie in a one-dimensional space. Burungale, Skinner, Tian, and Wan \cite{BSTW2024} proved the existence of $p$-adic zeta elements for elliptic curves over imaginary quadratic fields, establishing the Kobayashi conjecture for semistable curves at supersingular primes and presenting the first infinite families of non-CM elliptic curves satisfying the full BSD conjecture.
\end{remark}

\subsubsection*{Threshold axiom: higher Gross--Zagier}

\begin{axiom}[HGZ Bridge -- HGZ-BSD]\label{ax:hgz-bsd}
For an elliptic curve $E/\mathbb{Q}$ of analytic rank $r = \mathrm{ord}_{s=1} L(E,s)$, there exists:
\begin{enumerate}
\item[(HGZ-A)] \textbf{Canonical cycles:} A family of arithmetic cycles $\mathcal{Z}_r^{(i)}$ ($i = 1, \ldots, r$) in the appropriate Chow group or motivic cohomology, functorial under Hecke operators and base change.
\item[(HGZ-B)] \textbf{Height compatibility:} A canonical height pairing $\langle \cdot, \cdot \rangle_{\mathrm{ht}}$ on these cycles, compatible with the N\'eron--Tate height for $r=1$ and satisfying $\det(\langle \mathcal{Z}_r^{(i)}, \mathcal{Z}_r^{(j)} \rangle_{\mathrm{ht}}) = R_E \cdot (\text{local factors})$.
\item[(HGZ-C)] \textbf{Selmer control:} An Euler system of rank $r$ controlling the Selmer group $\mathrm{Sel}(E/\mathbb{Q})$, yielding $|\mathrm{Sha}(E/\mathbb{Q})| < \infty$ and the Selmer-side inequality $\mathrm{rank} \leq \mathrm{ord}$. The complementary inequality $\mathrm{ord} \leq \mathrm{rank}$ must come from the non-degenerate height/cycle part (HGZ-A/B), not from Selmer control alone.
\end{enumerate}
The conjunction (HGZ-A)+(HGZ-B)+(HGZ-C) implies the full BSD formula for rank $r$.
\end{axiom}

\begin{remark}[Module status]\label{rem:hgz-modules}
\begin{center}
\footnotesize
\resizebox{\columnwidth}{!}{%
\begin{tabular}{@{}lll@{}}
\toprule
\textbf{Module} & \textbf{Rank 1} & \textbf{Rank $\geq 2$} \\
\midrule
(A) Cycles & Heegner points (GZ 1986) & Diagonal cycles, partial \\
(B) Heights & GZ formula (1986) & Open: higher GZ formula \\
(C) Selmer & Kolyvagin (1988) & Partial: GSp(4) Euler systems \\
\bottomrule
\end{tabular}
}
\end{center}
The red block for rank $\geq 2$ is primarily \textbf{Module C}: existing Euler systems (Loeffler--Zerbes for GSp(4), Burungale--Skinner--Tian--Wan \cite{BSTW2024} for imaginary quadratic fields) provide partial control but do not yet yield a complete Kolyvagin-type bound on Sha for individual curves of rank $\geq 2$.
\end{remark}

\begin{remark}[No-go mechanisms]\label{rem:hgz-nogo}
HGZ-BSD fails if:
\begin{enumerate}
\item \textbf{Non-canonical cycle choice:} The higher cycles $\mathcal{Z}_r^{(i)}$ depend on auxiliary data (quaternion algebras, test vectors) in a way that cannot be uniformised---this is a language problem, not a computational one.
\item \textbf{Degenerate heights:} The height pairing $\langle \cdot, \cdot \rangle_{\mathrm{ht}}$ has a non-trivial kernel for $r \geq 2$, collapsing the determinant. Height saturation (Remark~\ref{conj:height-saturation-identity}) guards against this.
\item \textbf{Euler system weakness:} The rank-$r$ Euler system has insufficient variability to control all $p$-primary components of Sha simultaneously.
\end{enumerate}
\end{remark}

\begin{remark}[Module C as arithmetic underdetermination]\label{rem:module-c-underdetermined}
Without rank-$r$ Euler system control (Module~C), the Higher Gross--Zagier formula remains \emph{arithmetically underdetermined}: heights can be large without implying Selmer structure. This is the precise content of the red block: the bridge lemma (HGZ-A)+(HGZ-B)+(HGZ-C) $\Rightarrow$ $L^{(r)} \sim \det(\text{Heights})$ requires all three modules, and Module~C is the only one without even a partial result for general rank $r \geq 2$ curves. Loeffler--Zerbes \cite{LoefflerZerbes2024} provide partial GSp(4) control, but a full Kolyvagin-derivative argument for $|\mathrm{Sha}| < \infty$ at rank $\geq 2$ remains open.
\end{remark}

\subsubsection{Arakelov and motivic perspectives}

\begin{remark}[Universal Arakelov regulator]
\label{rem:arakelov-regulator}
The BSD formula involves both $p$-adic data (Selmer groups, Iwasawa modules) and archimedean data (periods $\Omega_E$, heights $\hat{h}$). The Arakelov intersection theory on the regular model $\mathcal{E} \to \mathrm{Spec}(\mathbb{Z})$ provides a unified framework: the Arakelov regulator $\hat{R}_E$ combines the archimedean N\'eron-Tate height with $p$-adic local contributions into a single adelic height. A ``universal Arakelov regulator'' as a roof over both $p$-adic and archimedean regulators would provide the natural bridge for hypothesis (H1): one formula $L^{(r)}(E,1) = C_r \cdot \hat{R}_E$ valid at all places simultaneously.
\end{remark}

\begin{remark}[Arakelov-motivic volume formula]
\label{rem:arakelov-volume}
One can attempt to interpret $L^{(r)}(E,1)$ as the measure of an ``arithmetic volume'' associated to an Arakelov vector bundle on the arithmetic surface $\mathcal{E}$. In this picture, $R_E$ is the volume of the fundamental domain of the Mordell-Weil lattice (with respect to the N\'eron-Tate metric), $\Omega_E$ is the archimedean volume of the fibre, and $|\Sha|$ is a ``defect volume'' measuring the failure of local-global compatibility. The BSD formula then asserts that the arithmetic Euler characteristic -- computed via the $L$-function -- equals this motivic volume. Connecting this to the Tamagawa number conjecture \cite{BlochKato1990,FontainePerrinRiou1994} would provide a conceptual derivation of the BSD formula. Indeed, the Bloch-Kato conjecture already provides the rigorous motivic framework for the uniqueness question discussed in Proposition~\ref{thm:uniqueness}: it defines Tamagawa numbers for general motives and derives the BSD formula as the special case for $h^1(E)(1)$, removing the rational ambiguity of the Beilinson conjectures.
\end{remark}

\begin{remark}[Hodge-Arakelov volumes: $\Sha$ as topological volume]
\label{rem:hodge-arakelov-sha}
The Hodge-Arakelov theory of Mochizuki interprets the comparison between de Rham and \'etale cohomology of elliptic curves as a ``discrete analogue'' of the Hodge decomposition. In this framework, $|\Sha|$ could be viewed as the measure of a topological volume -- the ``obstruction to global triviality'' -- that is forced to be finite by positivity constraints analogous to the Hodge-Riemann relations. If a positivity framework could be applied to this topological volume, finiteness of $\Sha$ (Axiom~I / hypothesis H3) might follow from positivity alone.
\end{remark}

\subsubsection{Algebraic and categorical approaches}

\begin{remark}[Higher Chow groups on Picard modular surfaces]
\label{rem:chow-groups}
For rank~$r \geq 2$, the relevant motives are no longer one-dimensional: the height pairing involves an $r \times r$ regulator matrix. The natural algebraic framework is provided by higher Chow groups $\mathrm{CH}^r(X, r)$ on Picard modular surfaces (Shimura varieties for $\mathrm{U}(2,1)$) or on self-products of the modular curve. The derived intersection pairing on these higher Chow groups generalizes the N\'eron-Tate pairing and could provide the $r \times r$ height matrix needed for hypothesis~(H1). The connection to $L$-functions is via the Beilinson conjectures, which predict that special values of $L$-functions are related to regulators on higher Chow groups.
\end{remark}

\begin{remark}[Derived Hecke algebras and Taylor-Wiles defect]
\label{rem:derived-hecke}
Venkatesh's derived Hecke algebra program relates the Taylor-Wiles method (patching completed cohomology) to derived algebraic geometry. The ``Taylor-Wiles defect'' -- the discrepancy between the dimension of a deformation ring and the dimension of the cohomology -- is conjectured to equal the rank of the Mordell-Weil group. If this could be made precise for elliptic curves, it would provide a new route to the rank part of BSD: $\rank E(\mathbb{Q}) = \text{TW defect} = \ord_{s=1} L(E,s)$.
\end{remark}

\subsubsection{Analytic and structural approaches}

\begin{remark}[$L$-function decomposition into Galois sectors]
\label{rem:galois-sectors}
The $L$-function $L(E,s)$ decomposes under twists by characters $\chi$ of $\Gal(\bar{\mathbb{Q}}/\mathbb{Q})$: the Hasse-Weil-Artin $L$-function $L(E, \chi, s)$ encodes the arithmetic of $E$ in the ``$\chi$-sector.'' By analogy with the topological sector decomposition in Yang-Mills theory (cf.\ the companion YM paper), one could decompose the BSD problem sector-by-sector: proving BSD for each twisted $L$-function $L(E, \chi, s)$ separately, then assembling the full result. Darmon and Rotger~\cite{DarmonRotger2017} pursue exactly this strategy for Artin representations, establishing BSD for Hasse-Weil-Artin $L$-functions in rank~1 cases.
\end{remark}

\begin{remark}[Cohomological sweeping processes for $\Sha$ finiteness]
\label{rem:sweeping-sha}
The finiteness of $\Sha$ (Axiom~I / hypothesis H3) can be reformulated as a statement about the global-to-local map in Galois cohomology: $\Sha(E/\mathbb{Q}) = \ker\!\left(H^1(\mathbb{Q}, E) \to \prod_v H^1(\mathbb{Q}_v, E)\right)$. A ``sweeping process'' approach from convex analysis would attempt to show that this kernel is finite by iteratively ``sweeping'' local obstructions: at each prime $v$, the local contribution $H^1(\mathbb{Q}_v, E)$ absorbs a definite proportion of the global classes, leaving only finitely many in the kernel. The key input would be a uniform bound on the ``absorption rate'' across primes, which is related to the Tamagawa numbers $c_v$ and the Lang-Trotter heuristics.
\end{remark}

\begin{remark}[O-minimal structures and $\Sha$ finiteness]
\label{rem:o-minimal-sha}
The o-minimal methods of Bakker-Brunebarbe-Tsimerman~\cite{BBT2023} (successfully applied to the Hodge conjecture, cf.\ the companion Hodge paper) could potentially be applied to $\Sha$ finiteness. The key idea: the period domain for the family of elliptic curves over $\mathbb{Q}$ is a definable analytic space in an o-minimal structure, and the set of curves with infinite $\Sha$ (if any exist) would form a definable subset. The tameness properties of o-minimal structures (finite stratification, controlled topology) could impose constraints on this locus, potentially showing it is empty. This is speculative but connects the BSD and Hodge programs at a structural level.
\end{remark}

\subsection{Numerical verification}
\label{sec:numerical}

The BSD formula has been verified to high precision for thousands of elliptic curves up to rank~4 (LMFDB~\cite{LMFDB}). The following computations are conditional diagnostics: they check consistency with BSD and with the positivity normal-form picture, but they do not establish (H1), (H2), or (H3):
\begin{itemize}
\item Under BSD, for every curve with $\rank E(\mathbb{Q}) = r$, the ratio $L^{(r)}(E,1)/(r! \cdot \Omega_E \cdot R_E \cdot \prod c_p)$ is expected to equal $|\Sha|/|E_{\mathrm{tors}}|^2$ -- a positive rational number, where $|\Sha|$ is itself a perfect square when $\Sha$ is finite and the Cassels--Tate pairing is alternating.
\item The regulator $R_E$ -- being the Gram determinant of a positive-definite form -- gives a conditional consistency scale for $|L^{(r)}(E,1)|$ once the remaining BSD factors are fixed. It is not an unconditional lower bound for the analytic derivative.
\item For rank-2 CM curves, the Selmer splitting predicts a factorization of the Kolyvagin bound into eigenspace contributions.
\end{itemize}

The accompanying Examples~\ref{ex:11a}--\ref{ex:389a} verify these predictions for specific curves of rank~0, 1, and~2.

\subsection{Post-Zookeeper transfer: a higher-rank module programme}
\label{subsec:zookeeper-bsd-transfer}

The post-Zookeeper use of the present paper is not to claim a new proof of
BSD in higher rank.  Rather, it turns the higher-rank obstruction into a
modular transfer problem.  The relevant five slots are:

In particular, the CORE closure of Riemann-$C2/MS2$ has no automatic BSD
consequence.  It only fixes the shared normal-form discipline: identify an
operator or form, an explicit defect, a finite approximation scheme, a
gap, and a limiting passage.  For BSD the red block remains higher-rank
Selmer control and the higher Gross--Zagier bridge, especially Module~C in
Remark~\ref{rem:hgz-modules}.

\begin{description}
  \item[Operator/form.] The Neron--Tate and $p$-adic height pairings on
  Mordell--Weil and Selmer directions.
  \item[Defect.] The BSD defect
  \[
    \delta_{\mathrm{BSD}}(E)
      := \operatorname{ord}_{s=1}L(E,s)-\operatorname{rank}E(\mathbb{Q}),
  \]
  together with the residual between the predicted leading term and the
  regulator expression.
  \item[Approximation.] Finite conductor, fixed rank, finite Selmer level,
  and finite layers in an Iwasawa tower.
  \item[Gap.] Non-degeneracy of the archimedean or $p$-adic regulator.
  \item[Limit.] The passage from $p$-adic leading-term information and
  Selmer structure to the archimedean BSD leading-term formula.
\end{description}

This suggests a conservative next target.  Instead of formulating a
general rank-$r$ proof, one should first isolate a rank-two transfer
statement of the form
\[
\begin{gathered}
  \text{higher Gross--Zagier input}\\
  +\ \text{Selmer control}\\
  +\ \text{$p$-adic height non-degeneracy}\\
  \Longrightarrow \text{rank-two height saturation}.
\end{gathered}
\]
Kim's higher Gross--Zagier and Selmer/Iwasawa work, together with the
Castella--et al. $p$-adic BSD framework, provide the natural literature
backbone for this programme.  The expected gain is diagnostic rather than
final: higher rank is decomposed into four separately testable modules
instead of remaining a single monolithic gap.

\begin{remark}[Determinant-line diagnostic and advice guardrail]
\label{rem:det-line-diagnostic}
The current working ledger sharpens the rank-two programme by treating the
regulator as a norm on a determinant line rather than as a post-hoc scalar.
For a candidate channel package $Z=(Z_1,\ldots,Z_r)$, set
\[
H_Z=(\langle Z_i,Z_j\rangle_{\mathrm{NT}})_{i,j},\qquad
G_r(E,Z)=\frac{\det H_Z}{\prod_i (H_Z)_{ii}},
\]
with the complementary degeneracy proxy
$\widetilde G_r(E,Z)=\lambda_{\min}(H_Z)/\mathrm{tr}(H_Z)$ when the entries
are available. The proposed defect ledger is
\[
\begin{aligned}
D_{\det}(E,Z)=(&d_{\mathrm{analytic\ line}},d_{\mathrm{compare}},d_{\mathrm{gap}},\\
&d_{\mathrm{selmer\ tail}},d_{\mathrm{advice}}).
\end{aligned}
\]
The final component records whether the channel choice came from analytic,
$p$-adic or Hecke data, or whether it used a known Mordell--Weil basis or
manual rank information. The latter is useful as a reference normalisation
but cannot count as bridge evidence. Thus leading-term, regulator and
Mordell--Weil information must not serve simultaneously as the target and as
the predictor; otherwise the diagnostic merely recycles the BSD quantity it
is supposed to test.

A rank-two signal is not counted unless the same pipeline fails on
pre-registered negative controls: duplicated channels $Z_2=Z_1$, rank-one
curves embedded as rank-two candidates, and conductor-matched controls with
channel rules fixed before the analytic comparison.
\end{remark}

\subsection*{Result classification}

This paper establishes a \textbf{structural normal-form characterisation within the BSD/Bloch--Kato motivic product class}: under four positivity axioms plus the explicit bridge hypotheses (H1)--(H3), the BSD formula is the expected admissible canonical-product shape for the leading term of $L^{(r)}(E,1)$. It is not a uniqueness claim among arbitrary functions. For rank~$\leq 1$, this is a genuine reformulation of proven results; for rank~$\geq 2$, it is a conditional structural analysis.

\begin{center}
\footnotesize
\resizebox{\columnwidth}{!}{%
\begin{tabular}{@{}ll@{}}
\toprule
\textbf{Result} & \textbf{Type} \\
\midrule
\multicolumn{2}{@{}l}{\textit{Proven (unconditional):}} \\
\midrule
N\'eron--Tate positivity (Axiom II) & Classical theorem \\
Rank $\leq 1$ BSD consequences & Theorem; normal-form reading interpretive \\
\midrule
\multicolumn{2}{@{}l}{\textit{Numerical evidence (not proof):}} \\
\midrule
Rank-2 consistency diagnostics & Conditional checks; advice audit pending \\
Regulator scaling $R \sim N^{0.10}$ & Exploratory regression \\
\midrule
\multicolumn{2}{@{}l}{\textit{Structural contributions:}} \\
\midrule
Axiom system (I--IV) & Definition / structural constraints \\
Uniqueness of BSD product shape & Product-class claim; conditional H1--H3 \\
Gross--Zagier as prototype & Interpretation \\
\midrule
\multicolumn{2}{@{}l}{\textit{Open (fundamental):}} \\
\midrule
Higher Gross--Zagier ($r \geq 2$) & Open \\
Higher Kolyvagin systems & Open \\
Sha finiteness ($r \geq 2$) & Open \\
\bottomrule
\end{tabular}
}
\end{center}

% ===================================================================
% 11. CONCLUSION
% ===================================================================
\section{Conclusion}
\label{sec:conclusion}

The Birch and Swinnerton-Dyer conjecture can be read, in this structural lens, as a discrepancy-control problem between analytic and geometric structures. The N\'eron-Tate height provides the canonical positivity structure: it is a positive-definite form on rational points that detects arithmetic directions via their energy cost.

The BSD formula is the expected canonical product normal form compatible with four axiomatic constraints plus the required bridge hypotheses: finite arithmetic capacity ($\Sha$ finiteness), cooperative height positivity (N\'eron-Tate), controlled Iwasawa growth, Euler-system composition coherence, height coupling and Selmer control. This does not reduce BSD to positivity alone; in higher rank it repackages BSD-level inputs into named bridge modules.

The rank $\leq 1$ proof is already compatible with this positivity No-Go reading: Gross-Zagier couples an analytic derivative to a non-negative height, and Kolyvagin's Euler system supplies Selmer control. For higher rank, the framework identifies two missing ingredients: non-degenerate higher Gross-Zagier height/cycle coupling and full higher-rank Selmer/Kolyvagin/Sha control. Height Saturation names the target but does not close either bridge.

For rank $\leq 1$, the contribution is conceptual: a normal-form reading of known theorems. For rank $\geq 2$, the contribution is a guarded research programme and diagnostic ledger for separating theorem input, conjectural bridge input, and advice-normalised numerical checks.

\begin{quote}
\textit{``BSD does not ask us to find rational points. It asks us to show that the $L$-function has no room for zeros without them.''}
\end{quote}


% ===================================================================
% ACKNOWLEDGMENTS
% ===================================================================
\begin{acknowledgments}
This work was conducted as independent research without institutional funding.

\textit{AI tools.} AI-based tools (Claude by Anthropic) were used for mathematical formalization and text generation. All mathematical content, conjectures, and responsibility for correctness lie solely with the author.

\textit{Funding information.} This research received no external funding.
\end{acknowledgments}


% ===================================================================
% BIBLIOGRAPHY
% ===================================================================
\begin{thebibliography}{35}
\setlength{\itemsep}{0pt}
\setlength{\parskip}{0pt}

\bibitem{BSD1965}
B.~J.~Birch \& H.~P.~F.~Swinnerton-Dyer (1965).
Notes on elliptic curves. II.
\textit{Journal f\"ur die reine und angewandte Mathematik}, 218, 79--108.
DOI: 10.1515/crll.1965.218.79.

\bibitem{GrossZagier1986}
B.~Gross \& D.~Zagier (1986).
Heegner points and derivatives of $L$-series.
\textit{Inventiones Mathematicae}, 84, 225--320.
DOI: 10.1007/BF01388809.

\bibitem{Kolyvagin1990}
V.~A.~Kolyvagin (1990).
Euler systems.
In \textit{The Grothendieck Festschrift}, Vol.~II, pp.\ 435--483. Birkh\"auser.
DOI: 10.1007/978-0-8176-4575-5\_11.

\bibitem{Kato2004}
K.~Kato (2004).
$p$-adic Hodge theory and values of zeta functions of modular forms.
\textit{Ast\'erisque}, 295, 117--290.

\bibitem{Kim2024}
C.-H.~Kim (2024).
A higher Gross-Zagier formula and the structure of Selmer groups.
To appear in \textit{Transactions of the AMS}.
arXiv:2203.12161.

\bibitem{SkinnerUrban2014}
C.~Skinner \& E.~Urban (2014).
The Iwasawa main conjectures for $\mathrm{GL}_2$.
\textit{Inventiones Mathematicae}, 195(1), 1--277.
DOI: 10.1007/s00222-013-0448-1.

\bibitem{MazurRubin2004}
B.~Mazur \& K.~Rubin (2004).
\textit{Kolyvagin Systems}.
Memoirs of the AMS, 168(799).
DOI: 10.1090/memo/0799.

\bibitem{DarmonRotger2017}
H.~Darmon \& V.~Rotger (2017).
Diagonal cycles and Euler systems II: The Birch and Swinnerton-Dyer conjecture for Hasse-Weil-Artin $L$-functions.
\textit{Journal of the AMS}, 30(3), 601--672.
DOI: 10.1090/jams/861.

\bibitem{BDP2013}
M.~Bertolini, H.~Darmon, \& K.~Prasanna (2013).
Generalized Heegner cycles and $p$-adic Rankin $L$-series.
\textit{Duke Mathematical Journal}, 162(6), 1033--1148.
DOI: 10.1215/00127094-2142056.

\bibitem{BurungaleCastellaKim2021}
A.~Burungale, F.~Castella, \& C.-H.~Kim (2021).
A proof of Perrin-Riou's Heegner point main conjecture.
\textit{Algebra \& Number Theory}, 15(7), 1627--1653.
DOI: 10.2140/ant.2021.15.1627.

\bibitem{BurnsSano2021}
D.~Burns \& T.~Sano (2021).
On the theory of higher rank Euler, Kolyvagin and Stark systems.
\textit{International Mathematics Research Notices}, 2021(13), 10118--10206.
DOI: 10.1093/imrn/rnz103.

\bibitem{BurnsSakamotoSano2025}
D.~Burns, R.~Sakamoto, \& T.~Sano (2025).
On the theory of higher rank Euler, Kolyvagin and Stark systems, II: the general theory.
To appear in \textit{Algebra \& Number Theory}.
arXiv:1805.08448.

\bibitem{LLZ2014}
A.~Lei, D.~Loeffler, \& S.~L.~Zerbes (2014).
Euler systems for Rankin-Selberg convolutions of modular forms.
\textit{Annals of Mathematics}, 180(2), 653--771.
DOI: 10.4007/annals.2014.180.2.6.

\bibitem{LSZ2020}
D.~Loeffler, C.~Skinner, \& S.~L.~Zerbes (2022).
Euler systems for $\mathrm{GSp}(4)$.
\textit{Journal of the European Mathematical Society}, 24(2), 669--733.
DOI: 10.4171/JEMS/1124.

\bibitem{Silverman2009}
J.~H.~Silverman (2009).
\textit{The Arithmetic of Elliptic Curves}, 2nd ed.
Graduate Texts in Mathematics 106, Springer.
DOI: 10.1007/978-0-387-09494-6.

\bibitem{Wiles1995}
A.~Wiles (1995).
Modular elliptic curves and Fermat's Last Theorem.
\textit{Annals of Mathematics}, 141(3), 443--551.
DOI: 10.2307/2118559.

\bibitem{Geiger2026e}
L.~Geiger (2026).
The Saturation Theorem: Axiomatic Constraints on Quantum Gravity from Cosmological Relaxation.
Zenodo preprint.
\href{https://doi.org/10.5281/zenodo.19036188}{doi:10.5281/zenodo.19036188}.

\bibitem{FST-Unified2026}
L.~Geiger (2026).
FST Spectrum Duality: The Renormalized Free Energy Principle as Physical Instantiation of Functional Stability Theory.
Zenodo preprint.
\href{https://doi.org/10.5281/zenodo.19036190}{doi:10.5281/zenodo.19036190}.

\bibitem{GeigerRH2026c}
L.~Geiger (2026).
From Landscape to Proof: The Even-Dominance Route to the Riemann Hypothesis.
Zenodo preprint.
\href{https://doi.org/10.5281/zenodo.19035640}{doi:10.5281/zenodo.19035640}.

\bibitem{GreenbergStevens1993}
R.~Greenberg \& G.~Stevens (1993).
$p$-adic $L$-functions and $p$-adic periods of modular forms.
\textit{Inventiones Mathematicae}, 111(2), 407--447.
DOI: 10.1007/BF01231294.

\bibitem{BhargavaShankar2015}
M.~Bhargava \& A.~Shankar (2015).
Ternary cubic forms having bounded invariants, and the existence of a positive proportion of elliptic curves having rank~0.
\textit{Annals of Mathematics}, 181(2), 587--621.
DOI: 10.4007/annals.2015.181.2.4.

\bibitem{LMFDB}
The LMFDB Collaboration (2024).
\textit{The L-functions and Modular Forms Database}.
\url{https://www.lmfdb.org}.

\bibitem{Nekovar2006}
J.~Nekov\'a\v{r} (2006).
Selmer complexes.
\textit{Ast\'erisque}, 310, vi+559.

\bibitem{Dokchitser2010}
T.~Dokchitser \& V.~Dokchitser (2010).
On the Birch-Swinnerton-Dyer quotients modulo squares.
\textit{Annals of Mathematics}, 172(1), 567--596.
DOI: 10.4007/annals.2010.172.567.

\bibitem{YZZ2013}
X.~Yuan, S.-W.~Zhang, \& W.~Zhang (2013).
\textit{The Gross-Zagier Formula on Shimura Curves}.
Annals of Mathematics Studies 184, Princeton University Press.

\bibitem{PerrinRiou1987}
B.~Perrin-Riou (1987).
Points de Heegner et d\'eriv\'ees de fonctions $L$ $p$-adiques.
\textit{Inventiones Mathematicae}, 89(3), 455--510.
DOI: 10.1007/BF01388982.

\bibitem{KLZ2017}
G.~Kings, D.~Loeffler, \& S.~L.~Zerbes (2017).
Rankin-Eisenstein classes and explicit reciprocity laws.
\textit{Cambridge Journal of Mathematics}, 5(1), 1--122.
DOI: 10.4310/CJM.2017.v5.n1.a1.

\bibitem{BBT2023}
B.~Bakker, B.~Brunebarbe \& J.~Tsimerman (2023).
o-minimal GAGA and a conjecture of Griffiths.
\textit{Inventiones Mathematicae}, 232(1), 163--228.
DOI: 10.1007/s00222-022-01166-1.

\bibitem{Smith2022}
A.~Smith, \textit{$2^\infty$-Selmer groups, $2^\infty$-class groups, and Goldfeld's conjecture}, arXiv:1702.02325v3, 2022.

\bibitem{LoefflerZerbes2024} D.~Loeffler and S.~L.~Zerbes, \textit{A universal Euler system for GSp(4)}, arXiv:2411.12576, 2024.

\bibitem{BSTW2024} A.~Burungale, C.~Skinner, Y.~Tian, and X.~Wan, \textit{Zeta elements for elliptic curves and applications}, arXiv:2409.01350, 2024.

\bibitem{Cesnavicius2018}
K.~\v{C}esnavi\v{c}ius (2018).
The Manin constant in the semistable case.
\textit{Compositio Mathematica}, 154(9), 1889--1920.
DOI: 10.1112/S0010437X18007273.

\bibitem{BlochKato1990}
S.~Bloch \& K.~Kato (1990).
$L$-functions and Tamagawa numbers of motives.
In \textit{The Grothendieck Festschrift}, Vol.~I, pp.\ 333--400. Birkh\"auser.
DOI: 10.1007/978-0-8176-4574-8\_9.

\bibitem{FontainePerrinRiou1994}
J.-M.~Fontaine \& B.~Perrin-Riou (1994).
Autour des conjectures de Bloch et Kato: cohomologie galoisienne et valeurs de fonctions~$L$.
In \textit{Motives} (Seattle 1991), Proc.\ Symp.\ Pure Math.\ 55, Part~1, pp.\ 599--706. AMS.

\bibitem{Szpiro1981}
L.~Szpiro (1981).
Propri\'et\'es num\'eriques du faisceau dualisant r\'elatif.
In: \textit{S\'eminaire sur les pinceaux de courbes de genre au moins deux},
Ast\'erisque \textbf{86}, 44--78.

\bibitem{GeigerIUT2026}
L.~Geiger (2026).
\textit{The Gap in IUTchIII, Corollary~3.12: From Attempted Repair to Structural Diagnosis (DRAFT)}.
Zenodo preprint.
\href{https://doi.org/10.5281/zenodo.20223276}{doi:10.5281/zenodo.20223276}.

\end{thebibliography}

\end{document}
